\documentclass[conference]{IEEEtran}
\usepackage{amsmath,amssymb}
\usepackage{graphicx}
\usepackage{booktabs}
\usepackage{xcolor}
\usepackage{tikz}
\usetikzlibrary{positioning,arrows.meta,fit,calc}
\usepackage{listings}
\lstset{basicstyle=\ttfamily\scriptsize,columns=fullflexible,frame=single,breaklines=true,aboveskip=4pt,belowskip=2pt}
\usepackage{url}
\usepackage[hidelinks]{hyperref}
\hypersetup{pdftitle={Can Agents Trust a Replication Verifier? Epistemic Contracts for Scientific Agent Tools},
  pdfauthor={Luxin Zhang}, pdfkeywords={agentic AI, epistemic contracts, verification, LLM-as-a-judge, MCP, replicability}}
% All result numbers in the prose are macros generated from results/ by src/make_numbers.py.
\newcommand{\aucBone}{0.627}
\newcommand{\sdBone}{0.017}
\newcommand{\ciBone}{[0.55, 0.71]}
\newcommand{\brierBone}{0.2409}
\newcommand{\bssBone}{2.6\%}
\newcommand{\aucMtwo}{0.620}
\newcommand{\sdMtwo}{0.020}
\newcommand{\ciMtwo}{[0.55, 0.70]}
\newcommand{\brierMtwo}{0.2406}
\newcommand{\bssMtwo}{2.8\%}
\newcommand{\aucDplan}{0.625}
\newcommand{\sdDplan}{0.019}
\newcommand{\ciDplan}{[0.54, 0.71]}
\newcommand{\brierDplan}{0.2429}
\newcommand{\bssDplan}{1.8\%}
\newcommand{\aucDach}{0.638}
\newcommand{\sdDach}{0.024}
\newcommand{\ciDach}{[0.58, 0.72]}
\newcommand{\brierDach}{0.2403}
\newcommand{\bssDach}{2.9\%}
\newcommand{\aucGonly}{0.591}
\newcommand{\sdGonly}{0.027}
\newcommand{\ciGonly}{[0.53, 0.69]}
\newcommand{\brierGonly}{0.2485}
\newcommand{\bssGonly}{-0.4\%}
\newcommand{\aucBtwo}{0.481}
\newcommand{\sdBtwo}{0.022}
\newcommand{\ciBtwo}{[0.37, 0.55]}
\newcommand{\brierBtwo}{0.2487}
\newcommand{\bssBtwo}{-0.5\%}
\newcommand{\aucBgbt}{0.620}
\newcommand{\sdBgbt}{0.020}
\newcommand{\ciBgbt}{[0.55, 0.72]}
\newcommand{\brierBgbt}{0.2483}
\newcommand{\bssBgbt}{-0.3\%}
\newcommand{\aucMtwolaunch}{0.620}
\newcommand{\sdMtwolaunch}{0.020}
\newcommand{\ciMtwolaunch}{[0.55, 0.70]}
\newcommand{\brierMtwolaunch}{0.2405}
\newcommand{\bssMtwolaunch}{2.8\%}
\newcommand{\aucBzero}{0.492}
\newcommand{\sdBzero}{0.001}
\newcommand{\ciBzero}{[0.37, 0.56]}
\newcommand{\brierBzero}{0.2474}
\newcommand{\bssBzero}{0.0\%}
\newcommand{\dBoneBzero}{\mbox{$+$0.156}}
\newcommand{\dciBoneBzero}{[\mbox{$+$0.025}, \mbox{$+$0.305}]}
\newcommand{\holmBoneBzero}{0.19}
\newcommand{\dbrierBoneBzero}{\mbox{$-$0.0087}}
\newcommand{\dbrierciBoneBzero}{[\mbox{$-$0.0215}, \mbox{$+$0.0036}]}
\newcommand{\dBfourBone}{\mbox{$-$0.013}}
\newcommand{\dciBfourBone}{[\mbox{$-$0.027}, \mbox{$-$0.001}]}
\newcommand{\holmBfourBone}{0.34}
\newcommand{\dbrierBfourBone}{\mbox{$+$0.0023}}
\newcommand{\dbrierciBfourBone}{[\mbox{$+$0.0004}, \mbox{$+$0.0044}]}
\newcommand{\dMoneBfour}{\mbox{$+$0.013}}
\newcommand{\dciMoneBfour}{[\mbox{$-$0.028}, \mbox{$+$0.046}]}
\newcommand{\holmMoneBfour}{1.00}
\newcommand{\dbrierMoneBfour}{\mbox{$-$0.0026}}
\newcommand{\dbrierciMoneBfour}{[\mbox{$-$0.0073}, \mbox{$+$0.0026}]}
\newcommand{\dMtwoMone}{\mbox{$+$0.000}}
\newcommand{\dciMtwoMone}{[\mbox{$-$0.017}, \mbox{$+$0.017}]}
\newcommand{\holmMtwoMone}{1.00}
\newcommand{\dbrierMtwoMone}{\mbox{$-$0.0002}}
\newcommand{\dbrierciMtwoMone}{[\mbox{$-$0.0027}, \mbox{$+$0.0025}]}
\newcommand{\dMtwoBone}{\mbox{$+$0.000}}
\newcommand{\dciMtwoBone}{[\mbox{$-$0.041}, \mbox{$+$0.033}]}
\newcommand{\holmMtwoBone}{1.00}
\newcommand{\dbrierMtwoBone}{\mbox{$-$0.0006}}
\newcommand{\dbrierciMtwoBone}{[\mbox{$-$0.0060}, \mbox{$+$0.0056}]}
\newcommand{\dMtwoNC}{\mbox{$+$0.015}}
\newcommand{\dciMtwoNC}{[\mbox{$-$0.018}, \mbox{$+$0.045}]}
\newcommand{\holmMtwoNC}{1.00}
\newcommand{\dbrierMtwoNC}{\mbox{$-$0.0021}}
\newcommand{\dbrierciMtwoNC}{[\mbox{$-$0.0068}, \mbox{$+$0.0030}]}
\newcommand{\dLeakMtwo}{\mbox{$-$0.010}}
\newcommand{\dciLeakMtwo}{[\mbox{$-$0.059}, \mbox{$+$0.035}]}
\newcommand{\holmLeakMtwo}{1.00}
\newcommand{\dbrierLeakMtwo}{\mbox{$+$0.0002}}
\newcommand{\dbrierciLeakMtwo}{[\mbox{$-$0.0077}, \mbox{$+$0.0089}]}
\newcommand{\dLeakG}{\mbox{$-$0.030}}
\newcommand{\dciLeakG}{[\mbox{$-$0.107}, \mbox{$+$0.050}]}
\newcommand{\holmLeakG}{1.00}
\newcommand{\dbrierLeakG}{\mbox{$-$0.0020}}
\newcommand{\dbrierciLeakG}{[\mbox{$-$0.0115}, \mbox{$+$0.0075}]}
\newcommand{\dDplanBone}{\mbox{$-$0.003}}
\newcommand{\dciDplanBone}{[\mbox{$-$0.037}, \mbox{$+$0.034}]}
\newcommand{\holmDplanBone}{1.00}
\newcommand{\dbrierDplanBone}{\mbox{$-$0.0003}}
\newcommand{\dbrierciDplanBone}{[\mbox{$-$0.0079}, \mbox{$+$0.0077}]}
\newcommand{\dDachDplan}{\mbox{$+$0.024}}
\newcommand{\dciDachDplan}{[\mbox{$-$0.005}, \mbox{$+$0.053}]}
\newcommand{\holmDachDplan}{1.00}
\newcommand{\dbrierDachDplan}{\mbox{$-$0.0028}}
\newcommand{\dbrierciDachDplan}{[\mbox{$-$0.0086}, \mbox{$+$0.0026}]}
\newcommand{\punBoneBzero}{0.02}
\newcommand{\selTopBone}{0.44}
\newcommand{\selTopCIBone}{[0.26, 0.68]}
\newcommand{\selHalfBone}{0.66}
\newcommand{\selAllBone}{0.63}
\newcommand{\selPosBone}{96\%}
\newcommand{\selTopMtwo}{0.67}
\newcommand{\selTopCIMtwo}{[0.38, 0.82]}
\newcommand{\selHalfMtwo}{0.66}
\newcommand{\selAllMtwo}{0.59}
\newcommand{\selPosMtwo}{89\%}
\newcommand{\selTopDplan}{0.48}
\newcommand{\selTopCIDplan}{[0.33, 0.70]}
\newcommand{\selHalfDplan}{0.65}
\newcommand{\selAllDplan}{0.64}
\newcommand{\selPosDplan}{89\%}
\newcommand{\selPosMin}{89\%}
\newcommand{\selPosMax}{96\%}
\newcommand{\permK}{30}
\newcommand{\permSeeds}{5}
\newcommand{\permReal}{0.632}
\newcommand{\permNullMean}{0.605}
\newcommand{\permNullLo}{0.580}
\newcommand{\permNullHi}{0.640}
\newcommand{\permP}{0.10}
\newcommand{\permBone}{0.639}
\newcommand{\permPbrier}{0.10}
\newcommand{\nCiting}{23{,}666}
\newcommand{\nCitingStrict}{8{,}268}
\newcommand{\earliestStart}{2019-08-06}
\newcommand{\nLaunch}{259}
\newcommand{\nPlanned}{219}
\newcommand{\plannedRatio}{1.33}
\newcommand{\logCitesStrict}{3.30}
\newcommand{\logCitesLeaky}{4.50}
\newcommand{\hasEvStrict}{97.8\%}
\newcommand{\hasEvLeaky}{100\%}
\newcommand{\fpRep}{8.9}
\newcommand{\fpNon}{4.7}
\newcommand{\fpAuc}{0.61}
\newcommand{\fpowRep}{0.55}
\newcommand{\fpowNon}{0.44}
\newcommand{\fpowAuc}{0.60}
\newcommand{\fmargRep}{16\%}
\newcommand{\fmargNon}{27\%}
\newcommand{\fexactRep}{67\%}
\newcommand{\fexactNon}{85\%}
\newcommand{\fxfRep}{0.32}
\newcommand{\fxfNon}{0.42}
\newcommand{\fxfAuc}{0.62}
\newcommand{\fcitesAuc}{0.52}
\newcommand{\subPsych}{0.53}
\newcommand{\subPsychN}{94}
\newcommand{\subEcon}{0.78}
\newcommand{\subEconN}{38}
\newcommand{\subPol}{0.74}
\newcommand{\subPolN}{45}
\newcommand{\starterAuc}{0.523}
\newcommand{\starterBrier}{0.267}
\newcommand{\starterPrevBrier}{0.247}
\newcommand{\sensLo}{0.60}
\newcommand{\sensHi}{0.64}
\newcommand{\nTests}{9}
\newcommand{\nMutants}{5}
\newcommand{\nKilled}{5}
\newcommand{\locPython}{2877}
\newcommand{\pwBone}{0.632}
\newcommand{\pwMtwo}{0.602}
\newcommand{\pwDplan}{0.660}
\newcommand{\pwBase}{49.3\%}
\newcommand{\pwBssBone}{4.1\%}
\newcommand{\maxClaimsPaper}{18}
\newcommand{\llmMiniIdAuc}{0.600}
\newcommand{\llmMiniIdCi}{[0.52, 0.68]}
\newcommand{\llmMiniIdBrier}{0.2753}
\newcommand{\llmMiniIdVsBone}{\mbox{$-$0.028}}
\newcommand{\llmMiniIdVsBoneCi}{[\mbox{$-$0.106}, \mbox{$+$0.050}]}
\newcommand{\llmMiniIdTop}{0.70}
\newcommand{\llmMiniIdAll}{0.57}
\newcommand{\llmMiniIdSd}{0.19}
\newcommand{\llmMiniIdTopCi}{[0.48, 0.92]}
\newcommand{\llmMiniIdGainCi}{[\mbox{$-$0.08}, \mbox{$+$0.34}]}
\newcommand{\llmMiniIdTopPos}{93\%}
\newcommand{\llmMiniDeidAuc}{0.570}
\newcommand{\llmMiniDeidCi}{[0.48, 0.66]}
\newcommand{\llmMiniDeidBrier}{0.2758}
\newcommand{\llmMiniDeidVsBone}{\mbox{$-$0.058}}
\newcommand{\llmMiniDeidVsBoneCi}{[\mbox{$-$0.160}, \mbox{$+$0.038}]}
\newcommand{\llmMiniDeidTop}{0.56}
\newcommand{\llmMiniDeidAll}{0.57}
\newcommand{\llmMiniDeidSd}{0.19}
\newcommand{\llmMiniDeidTopCi}{[0.32, 0.81]}
\newcommand{\llmMiniDeidGainCi}{[\mbox{$-$0.23}, \mbox{$+$0.22}]}
\newcommand{\llmMiniDeidTopPos}{89\%}
\newcommand{\llmMiniIdDeid}{\mbox{$+$0.030}}
\newcommand{\llmMiniIdDeidCi}{[\mbox{$-$0.028}, \mbox{$+$0.098}]}
\newcommand{\llmMiniAwareReal}{42.7\%}
\newcommand{\llmMiniAwareDecoy}{24.0\%}
\newcommand{\llmMiniNamed}{77}
\newcommand{\llmMiniAwareRealCi}{[37, 49]\%}
\newcommand{\llmMiniAwareDecoyCi}{[17, 33]\%}
\newcommand{\llmMiniNamedDecoy}{12}
\newcommand{\llmMiniNamedAcc}{0.64}
\newcommand{\llmMiniNamedMaj}{0.56}
\newcommand{\llmMiniCost}{0.38}
\newcommand{\llmMiniFail}{0}
\newcommand{\llmMiniEffort}{none}
\newcommand{\llmSolIdAuc}{0.663}
\newcommand{\llmSolIdCi}{[0.58, 0.74]}
\newcommand{\llmSolIdBrier}{0.2594}
\newcommand{\llmSolIdVsBone}{\mbox{$+$0.036}}
\newcommand{\llmSolIdVsBoneCi}{[\mbox{$-$0.031}, \mbox{$+$0.105}]}
\newcommand{\llmSolIdTop}{0.78}
\newcommand{\llmSolIdAll}{0.61}
\newcommand{\llmSolIdSd}{0.18}
\newcommand{\llmSolIdTopCi}{[0.58, 0.97]}
\newcommand{\llmSolIdGainCi}{[\mbox{$-$0.00}, \mbox{$+$0.36}]}
\newcommand{\llmSolIdTopPos}{100\%}
\newcommand{\llmSolDeidAuc}{0.612}
\newcommand{\llmSolDeidCi}{[0.53, 0.70]}
\newcommand{\llmSolDeidBrier}{0.2627}
\newcommand{\llmSolDeidVsBone}{\mbox{$-$0.016}}
\newcommand{\llmSolDeidVsBoneCi}{[\mbox{$-$0.107}, \mbox{$+$0.073}]}
\newcommand{\llmSolDeidTop}{0.74}
\newcommand{\llmSolDeidAll}{0.58}
\newcommand{\llmSolDeidSd}{0.18}
\newcommand{\llmSolDeidTopCi}{[0.52, 0.93]}
\newcommand{\llmSolDeidGainCi}{[\mbox{$-$0.04}, \mbox{$+$0.35}]}
\newcommand{\llmSolDeidTopPos}{96\%}
\newcommand{\llmSolIdDeid}{\mbox{$+$0.051}}
\newcommand{\llmSolIdDeidCi}{[\mbox{$+$0.008}, \mbox{$+$0.101}]}
\newcommand{\llmSolAwareReal}{2.2\%}
\newcommand{\llmSolAwareDecoy}{2.0\%}
\newcommand{\llmSolNamed}{6}
\newcommand{\llmSolAwareRealCi}{[1, 5]\%}
\newcommand{\llmSolAwareDecoyCi}{[1, 7]\%}
\newcommand{\llmSolNamedDecoy}{2}
\newcommand{\llmSolNamedAcc}{1.00}
\newcommand{\llmSolNamedMaj}{0.83}
\newcommand{\llmSolCost}{1.33}
\newcommand{\llmSolFail}{0}
\newcommand{\llmSolEffort}{low}
\newcommand{\llmHaikuIdAuc}{0.566}
\newcommand{\llmHaikuIdCi}{[0.48, 0.65]}
\newcommand{\llmHaikuIdBrier}{0.2541}
\newcommand{\llmHaikuIdVsBone}{\mbox{$-$0.061}}
\newcommand{\llmHaikuIdVsBoneCi}{[\mbox{$-$0.154}, \mbox{$+$0.026}]}
\newcommand{\llmHaikuIdTop}{0.52}
\newcommand{\llmHaikuIdAll}{0.56}
\newcommand{\llmHaikuIdSd}{0.15}
\newcommand{\llmHaikuIdTopCi}{[0.39, 0.81]}
\newcommand{\llmHaikuIdGainCi}{[\mbox{$-$0.15}, \mbox{$+$0.22}]}
\newcommand{\llmHaikuIdTopPos}{63\%}
\newcommand{\llmHaikuDeidAuc}{0.592}
\newcommand{\llmHaikuDeidCi}{[0.51, 0.68]}
\newcommand{\llmHaikuDeidBrier}{0.2561}
\newcommand{\llmHaikuDeidVsBone}{\mbox{$-$0.036}}
\newcommand{\llmHaikuDeidVsBoneCi}{[\mbox{$-$0.132}, \mbox{$+$0.057}]}
\newcommand{\llmHaikuDeidTop}{0.70}
\newcommand{\llmHaikuDeidAll}{0.53}
\newcommand{\llmHaikuDeidSd}{0.15}
\newcommand{\llmHaikuDeidTopCi}{[0.36, 0.88]}
\newcommand{\llmHaikuDeidGainCi}{[\mbox{$-$0.15}, \mbox{$+$0.34}]}
\newcommand{\llmHaikuDeidTopPos}{81\%}
\newcommand{\llmHaikuIdDeid}{\mbox{$-$0.026}}
\newcommand{\llmHaikuIdDeidCi}{[\mbox{$-$0.093}, \mbox{$+$0.047}]}
\newcommand{\llmHaikuAwareReal}{0.0\%}
\newcommand{\llmHaikuAwareDecoy}{0.0\%}
\newcommand{\llmHaikuNamed}{0}
\newcommand{\llmHaikuAwareRealCi}{[0, 1]\%}
\newcommand{\llmHaikuAwareDecoyCi}{[0, 4]\%}
\newcommand{\llmHaikuNamedDecoy}{0}
\newcommand{\llmHaikuNamedAcc}{--}
\newcommand{\llmHaikuNamedMaj}{--}
\newcommand{\llmHaikuCost}{0.65}
\newcommand{\llmHaikuFail}{0}
\newcommand{\llmHaikuEffort}{default}
\newcommand{\llmOpusIdAuc}{0.655}
\newcommand{\llmOpusIdCi}{[0.57, 0.74]}
\newcommand{\llmOpusIdBrier}{0.2392}
\newcommand{\llmOpusIdVsBone}{\mbox{$+$0.027}}
\newcommand{\llmOpusIdVsBoneCi}{[\mbox{$-$0.044}, \mbox{$+$0.103}]}
\newcommand{\llmOpusIdTop}{0.89}
\newcommand{\llmOpusIdAll}{0.62}
\newcommand{\llmOpusIdSd}{0.23}
\newcommand{\llmOpusIdTopCi}{[0.67, 1.00]}
\newcommand{\llmOpusIdGainCi}{[\mbox{$+$0.07}, \mbox{$+$0.38}]}
\newcommand{\llmOpusIdTopPos}{89\%}
\newcommand{\llmOpusDeidAuc}{0.620}
\newcommand{\llmOpusDeidCi}{[0.53, 0.70]}
\newcommand{\llmOpusDeidBrier}{0.2478}
\newcommand{\llmOpusDeidVsBone}{\mbox{$-$0.007}}
\newcommand{\llmOpusDeidVsBoneCi}{[\mbox{$-$0.097}, \mbox{$+$0.077}]}
\newcommand{\llmOpusDeidTop}{0.70}
\newcommand{\llmOpusDeidAll}{0.60}
\newcommand{\llmOpusDeidSd}{0.20}
\newcommand{\llmOpusDeidTopCi}{[0.41, 0.93]}
\newcommand{\llmOpusDeidGainCi}{[\mbox{$-$0.16}, \mbox{$+$0.32}]}
\newcommand{\llmOpusDeidTopPos}{93\%}
\newcommand{\llmOpusIdDeid}{\mbox{$+$0.035}}
\newcommand{\llmOpusIdDeidCi}{[\mbox{$-$0.024}, \mbox{$+$0.092}]}
\newcommand{\llmOpusAwareReal}{0.0\%}
\newcommand{\llmOpusAwareDecoy}{0.0\%}
\newcommand{\llmOpusNamed}{3}
\newcommand{\llmOpusAwareRealCi}{[0, 1]\%}
\newcommand{\llmOpusAwareDecoyCi}{[0, 4]\%}
\newcommand{\llmOpusNamedDecoy}{0}
\newcommand{\llmOpusNamedAcc}{0.67}
\newcommand{\llmOpusNamedMaj}{0.67}
\newcommand{\llmOpusCost}{3.33}
\newcommand{\llmOpusFail}{0}
\newcommand{\llmOpusEffort}{default}
\newcommand{\llmBoneAuc}{0.628}
\newcommand{\llmN}{274}
\newcommand{\llmBoneBrier}{0.2387}
\newcommand{\llmBzeroBrier}{0.2474}
\newcommand{\llmBoneTop}{0.44}
\newcommand{\llmCalls}{3688}
\newcommand{\llmCostTotal}{5.68}
\newcommand{\rcMiniPlattBrier}{0.2438}
\newcommand{\rcMiniPlattAuc}{0.593}
\newcommand{\rcMiniPlattBss}{+1.5\%}
\newcommand{\rcMiniPlattVsBone}{\mbox{$+$0.0051}}
\newcommand{\rcMiniPlattVsBoneCi}{[\mbox{$-$0.0062}, \mbox{$+$0.0161}]}
\newcommand{\rcMiniPlattVsBzeroCi}{[\mbox{$-$0.0148}, \mbox{$+$0.0072}]}
\newcommand{\rcMiniIsoBrier}{0.2514}
\newcommand{\rcMiniIsoAuc}{0.545}
\newcommand{\rcMiniIsoBss}{-1.6\%}
\newcommand{\rcMiniIsoVsBone}{\mbox{$+$0.0127}}
\newcommand{\rcMiniIsoVsBoneCi}{[\mbox{$-$0.0019}, \mbox{$+$0.0275}]}
\newcommand{\rcMiniIsoVsBzeroCi}{[\mbox{$-$0.0126}, \mbox{$+$0.0210}]}
\newcommand{\rcSolPlattBrier}{0.2313}
\newcommand{\rcSolPlattAuc}{0.655}
\newcommand{\rcSolPlattBss}{+6.5\%}
\newcommand{\rcSolPlattVsBone}{\mbox{$-$0.0074}}
\newcommand{\rcSolPlattVsBoneCi}{[\mbox{$-$0.0229}, \mbox{$+$0.0074}]}
\newcommand{\rcSolPlattVsBzeroCi}{[\mbox{$-$0.0369}, \mbox{$+$0.0057}]}
\newcommand{\rcSolIsoBrier}{0.2398}
\newcommand{\rcSolIsoAuc}{0.623}
\newcommand{\rcSolIsoBss}{+3.1\%}
\newcommand{\rcSolIsoVsBone}{\mbox{$+$0.0011}}
\newcommand{\rcSolIsoVsBoneCi}{[\mbox{$-$0.0193}, \mbox{$+$0.0205}]}
\newcommand{\rcSolIsoVsBzeroCi}{[\mbox{$-$0.0335}, \mbox{$+$0.0186}]}
\newcommand{\rcHaikuPlattBrier}{0.2466}
\newcommand{\rcHaikuPlattAuc}{0.531}
\newcommand{\rcHaikuPlattBss}{+0.3\%}
\newcommand{\rcHaikuPlattVsBone}{\mbox{$+$0.0079}}
\newcommand{\rcHaikuPlattVsBoneCi}{[\mbox{$-$0.0047}, \mbox{$+$0.0195}]}
\newcommand{\rcHaikuPlattVsBzeroCi}{[\mbox{$-$0.0102}, \mbox{$+$0.0090}]}
\newcommand{\rcHaikuIsoBrier}{0.2515}
\newcommand{\rcHaikuIsoAuc}{0.473}
\newcommand{\rcHaikuIsoBss}{-1.7\%}
\newcommand{\rcHaikuIsoVsBone}{\mbox{$+$0.0129}}
\newcommand{\rcHaikuIsoVsBoneCi}{[\mbox{$-$0.0005}, \mbox{$+$0.0253}]}
\newcommand{\rcHaikuIsoVsBzeroCi}{[\mbox{$-$0.0074}, \mbox{$+$0.0155}]}
\newcommand{\rcOpusPlattBrier}{0.2323}
\newcommand{\rcOpusPlattAuc}{0.639}
\newcommand{\rcOpusPlattBss}{+6.1\%}
\newcommand{\rcOpusPlattVsBone}{\mbox{$-$0.0063}}
\newcommand{\rcOpusPlattVsBoneCi}{[\mbox{$-$0.0258}, \mbox{$+$0.0132}]}
\newcommand{\rcOpusPlattVsBzeroCi}{[\mbox{$-$0.0399}, \mbox{$+$0.0104}]}
\newcommand{\rcOpusIsoBrier}{0.2354}
\newcommand{\rcOpusIsoAuc}{0.593}
\newcommand{\rcOpusIsoBss}{+4.8\%}
\newcommand{\rcOpusIsoVsBone}{\mbox{$-$0.0032}}
\newcommand{\rcOpusIsoVsBoneCi}{[\mbox{$-$0.0225}, \mbox{$+$0.0166}]}
\newcommand{\rcOpusIsoVsBzeroCi}{[\mbox{$-$0.0368}, \mbox{$+$0.0145}]}
\newcommand{\mutTotal}{251}
\newcommand{\mutToolKilled}{134}
\newcommand{\mutToolScore}{53\%}
\newcommand{\mutVoneKilled}{153}
\newcommand{\mutVoneScore}{61\%}
\newcommand{\mutMissed}{19}
\newcommand{\mutVtwoKilled}{204}
\newcommand{\mutVtwoScore}{81\%}
\newcommand{\mutVtwoTimeout}{0}
\newcommand{\mutScore}{81\%}
\newcommand{\mutSurv}{42}
\newcommand{\mutVthreeKilled}{209}
\newcommand{\mutVthreeScore}{83\%}
\newcommand{\mutVthreeExEq}{86\%}
\newcommand{\mutScoreExEq}{84\%}
\newcommand{\mutClsEquivalent}{7}
\newcommand{\mutClsExplanationDrivers}{8}
\newcommand{\mutClsCliEntryPoint}{9}
\newcommand{\mutClsModelHyperparameters}{13}
\newcommand{\mutClsMessageOrRounding}{5}
\newcommand{\agMiniAUEscalate}{66\%}
\newcommand{\agMiniAUAuto}{34\%}
\newcommand{\agMiniAUUnsafeaccept}{1\%}
\newcommand{\agMiniAUWrongreject}{14\%}
\newcommand{\agMiniAUAutoerror}{15\%}
\newcommand{\agMiniAUAutoacc}{0.57}
\newcommand{\agMiniACEscalate}{93\%}
\newcommand{\agMiniACAuto}{7\%}
\newcommand{\agMiniACUnsafeaccept}{0\%}
\newcommand{\agMiniACWrongreject}{3\%}
\newcommand{\agMiniACAutoerror}{4\%}
\newcommand{\agMiniACAutoacc}{0.50}
\newcommand{\agMiniADEscalate}{\mbox{$+$26.3}}
\newcommand{\agMiniADEscalateCi}{[\mbox{$+$18}, \mbox{$+$34}]}
\newcommand{\agMiniADUnsafeaccept}{\mbox{$-$0.4}}
\newcommand{\agMiniADUnsafeacceptCi}{[\mbox{$-$2}, \mbox{$+$1}]}
\newcommand{\agMiniADAutoerror}{\mbox{$-$10.9}}
\newcommand{\agMiniADAutoerrorCi}{[\mbox{$-$16}, \mbox{$-$6}]}
\newcommand{\agMiniAN}{274}
\newcommand{\agMiniABreakeven}{0.42}
\newcommand{\agMiniACRej}{6\%}
\newcommand{\agMiniAConVsPla}{\mbox{$-$5.1}}
\newcommand{\agMiniAConVsPlaCi}{[\mbox{$-$9.2}, \mbox{$-$1.7}]}
\newcommand{\agMiniAUnpVsNo}{\mbox{$+$14.6}}
\newcommand{\agMiniAUnpVsNoCi}{[\mbox{$+$9.8}, \mbox{$+$19.8}]}
\newcommand{\agMiniAConVsNo}{\mbox{$+$3.6}}
\newcommand{\agMiniAConVsNoCi}{[\mbox{$+$1.5}, \mbox{$+$5.8}]}
\newcommand{\agMiniAUPtwoErr}{17.5\%}
\newcommand{\agMiniAUPtwoRej}{33\%}
\newcommand{\agMiniAPlaPtwoErr}{9.1\%}
\newcommand{\agMiniAPlaPtwoRej}{17\%}
\newcommand{\agMiniACPtwoErr}{1.5\%}
\newcommand{\agMiniACPtwoRej}{3\%}
\newcommand{\agMiniANotoolEsc}{100\%}
\newcommand{\agMiniANotoolErr}{0.0\%}
\newcommand{\agMiniANotoolRej}{0\%}
\newcommand{\agMiniAPlaceboEsc}{82\%}
\newcommand{\agMiniAPlaceboErr}{8.8\%}
\newcommand{\agMiniAPlaceboRej}{15\%}
\newcommand{\agMiniBUEscalate}{54\%}
\newcommand{\agMiniBUAuto}{46\%}
\newcommand{\agMiniBUUnsafeaccept}{1\%}
\newcommand{\agMiniBUWrongreject}{19\%}
\newcommand{\agMiniBUAutoerror}{20\%}
\newcommand{\agMiniBUAutoacc}{0.56}
\newcommand{\agMiniBCEscalate}{100\%}
\newcommand{\agMiniBCAuto}{0\%}
\newcommand{\agMiniBCUnsafeaccept}{0\%}
\newcommand{\agMiniBCWrongreject}{0\%}
\newcommand{\agMiniBCAutoerror}{0\%}
\newcommand{\agMiniBCAutoacc}{--}
\newcommand{\agMiniBDEscalate}{\mbox{$+$45.6}}
\newcommand{\agMiniBDEscalateCi}{[\mbox{$+$34}, \mbox{$+$57}]}
\newcommand{\agMiniBDUnsafeaccept}{\mbox{$-$1.2}}
\newcommand{\agMiniBDUnsafeacceptCi}{[\mbox{$-$3}, \mbox{$+$0}]}
\newcommand{\agMiniBDAutoerror}{\mbox{$-$20.1}}
\newcommand{\agMiniBDAutoerrorCi}{[\mbox{$-$29}, \mbox{$-$12}]}
\newcommand{\agMiniBN}{259}
\newcommand{\agMiniBBreakeven}{0.44}
\newcommand{\agMiniBCRej}{0\%}
\newcommand{\agSolAUEscalate}{81\%}
\newcommand{\agSolAUAuto}{19\%}
\newcommand{\agSolAUUnsafeaccept}{0\%}
\newcommand{\agSolAUWrongreject}{9\%}
\newcommand{\agSolAUAutoerror}{9\%}
\newcommand{\agSolAUAutoacc}{0.53}
\newcommand{\agSolACEscalate}{100\%}
\newcommand{\agSolACAuto}{0\%}
\newcommand{\agSolACUnsafeaccept}{0\%}
\newcommand{\agSolACWrongreject}{0\%}
\newcommand{\agSolACAutoerror}{0\%}
\newcommand{\agSolACAutoacc}{--}
\newcommand{\agSolADEscalate}{\mbox{$+$19.3}}
\newcommand{\agSolADEscalateCi}{[\mbox{$+$13}, \mbox{$+$26}]}
\newcommand{\agSolADUnsafeaccept}{\mbox{$+$0.0}}
\newcommand{\agSolADUnsafeacceptCi}{[\mbox{$+$0}, \mbox{$+$0}]}
\newcommand{\agSolADAutoerror}{\mbox{$-$9.1}}
\newcommand{\agSolADAutoerrorCi}{[\mbox{$-$14}, \mbox{$-$5}]}
\newcommand{\agSolAN}{274}
\newcommand{\agSolABreakeven}{0.47}
\newcommand{\agSolACRej}{0\%}
\newcommand{\agSolAConVsPla}{\mbox{$-$2.9}}
\newcommand{\agSolAConVsPlaCi}{[\mbox{$-$5.2}, \mbox{$-$1.1}]}
\newcommand{\agSolAUnpVsNo}{\mbox{$+$9.1}}
\newcommand{\agSolAUnpVsNoCi}{[\mbox{$+$5.2}, \mbox{$+$13.5}]}
\newcommand{\agSolAConVsNo}{\mbox{$+$0.0}}
\newcommand{\agSolAConVsNoCi}{[\mbox{$+$0.0}, \mbox{$+$0.0}]}
\newcommand{\agSolAUPtwoErr}{1.8\%}
\newcommand{\agSolAUPtwoRej}{4\%}
\newcommand{\agSolAPlaPtwoErr}{0.4\%}
\newcommand{\agSolAPlaPtwoRej}{0\%}
\newcommand{\agSolACPtwoErr}{0.0\%}
\newcommand{\agSolACPtwoRej}{0\%}
\newcommand{\agSolANotoolEsc}{100\%}
\newcommand{\agSolANotoolErr}{0.0\%}
\newcommand{\agSolANotoolRej}{0\%}
\newcommand{\agSolAPlaceboEsc}{95\%}
\newcommand{\agSolAPlaceboErr}{2.9\%}
\newcommand{\agSolAPlaceboRej}{5\%}
\newcommand{\agSolBUEscalate}{75\%}
\newcommand{\agSolBUAuto}{25\%}
\newcommand{\agSolBUUnsafeaccept}{1\%}
\newcommand{\agSolBUWrongreject}{11\%}
\newcommand{\agSolBUAutoerror}{12\%}
\newcommand{\agSolBUAutoacc}{0.52}
\newcommand{\agSolBCEscalate}{100\%}
\newcommand{\agSolBCAuto}{0\%}
\newcommand{\agSolBCUnsafeaccept}{0\%}
\newcommand{\agSolBCWrongreject}{0\%}
\newcommand{\agSolBCAutoerror}{0\%}
\newcommand{\agSolBCAutoacc}{--}
\newcommand{\agSolBDEscalate}{\mbox{$+$24.7}}
\newcommand{\agSolBDEscalateCi}{[\mbox{$+$15}, \mbox{$+$35}]}
\newcommand{\agSolBDUnsafeaccept}{\mbox{$-$0.8}}
\newcommand{\agSolBDUnsafeacceptCi}{[\mbox{$-$2}, \mbox{$+$0}]}
\newcommand{\agSolBDAutoerror}{\mbox{$-$12.0}}
\newcommand{\agSolBDAutoerrorCi}{[\mbox{$-$20}, \mbox{$-$6}]}
\newcommand{\agSolBN}{259}
\newcommand{\agSolBBreakeven}{0.48}
\newcommand{\agSolBCRej}{0\%}
\newcommand{\agHaikuAUEscalate}{59\%}
\newcommand{\agHaikuAUAuto}{41\%}
\newcommand{\agHaikuAUUnsafeaccept}{1\%}
\newcommand{\agHaikuAUWrongreject}{16\%}
\newcommand{\agHaikuAUAutoerror}{18\%}
\newcommand{\agHaikuAUAutoacc}{0.56}
\newcommand{\agHaikuACEscalate}{1\%}
\newcommand{\agHaikuACAuto}{99\%}
\newcommand{\agHaikuACUnsafeaccept}{1\%}
\newcommand{\agHaikuACWrongreject}{53\%}
\newcommand{\agHaikuACAutoerror}{54\%}
\newcommand{\agHaikuACAutoacc}{0.45}
\newcommand{\agHaikuADEscalate}{\mbox{$-$58.0}}
\newcommand{\agHaikuADEscalateCi}{[\mbox{$-$66}, \mbox{$-$50}]}
\newcommand{\agHaikuADUnsafeaccept}{\mbox{$-$0.4}}
\newcommand{\agHaikuADUnsafeacceptCi}{[\mbox{$-$2}, \mbox{$+$1}]}
\newcommand{\agHaikuADAutoerror}{\mbox{$+$36.1}}
\newcommand{\agHaikuADAutoerrorCi}{[\mbox{$+$27}, \mbox{$+$44}]}
\newcommand{\agHaikuAN}{274}
\newcommand{\agHaikuAVtwoEsc}{100\%}
\newcommand{\agHaikuAVtwoErr}{0\%}
\newcommand{\agHaikuACRej}{96\%}
\newcommand{\agHaikuAConVsPla}{\mbox{$+$38.3}}
\newcommand{\agHaikuAConVsPlaCi}{[\mbox{$+$30.2}, \mbox{$+$46.7}]}
\newcommand{\agHaikuAUnpVsNo}{\mbox{$+$17.5}}
\newcommand{\agHaikuAUnpVsNoCi}{[\mbox{$+$11.9}, \mbox{$+$23.4}]}
\newcommand{\agHaikuAConVsNo}{\mbox{$+$53.6}}
\newcommand{\agHaikuAConVsNoCi}{[\mbox{$+$45.7}, \mbox{$+$60.9}]}
\newcommand{\agHaikuAUPtwoErr}{31.0\%}
\newcommand{\agHaikuAUPtwoRej}{50\%}
\newcommand{\agHaikuAPlaPtwoErr}{33.2\%}
\newcommand{\agHaikuAPlaPtwoRej}{50\%}
\newcommand{\agHaikuACPtwoErr}{50.4\%}
\newcommand{\agHaikuACPtwoRej}{91\%}
\newcommand{\agHaikuANotoolEsc}{98\%}
\newcommand{\agHaikuANotoolErr}{0.4\%}
\newcommand{\agHaikuANotoolRej}{2\%}
\newcommand{\agHaikuAPlaceboEsc}{63\%}
\newcommand{\agHaikuAPlaceboErr}{15.7\%}
\newcommand{\agHaikuAPlaceboRej}{32\%}
\newcommand{\agHaikuBUEscalate}{49\%}
\newcommand{\agHaikuBUAuto}{51\%}
\newcommand{\agHaikuBUUnsafeaccept}{2\%}
\newcommand{\agHaikuBUWrongreject}{20\%}
\newcommand{\agHaikuBUAutoerror}{22\%}
\newcommand{\agHaikuBUAutoacc}{0.57}
\newcommand{\agHaikuBCEscalate}{85\%}
\newcommand{\agHaikuBCAuto}{15\%}
\newcommand{\agHaikuBCUnsafeaccept}{0\%}
\newcommand{\agHaikuBCWrongreject}{7\%}
\newcommand{\agHaikuBCAutoerror}{7\%}
\newcommand{\agHaikuBCAutoacc}{0.53}
\newcommand{\agHaikuBDEscalate}{\mbox{$+$35.9}}
\newcommand{\agHaikuBDEscalateCi}{[\mbox{$+$24}, \mbox{$+$47}]}
\newcommand{\agHaikuBDUnsafeaccept}{\mbox{$-$1.9}}
\newcommand{\agHaikuBDUnsafeacceptCi}{[\mbox{$-$4}, \mbox{$-$0}]}
\newcommand{\agHaikuBDAutoerror}{\mbox{$-$14.7}}
\newcommand{\agHaikuBDAutoerrorCi}{[\mbox{$-$22}, \mbox{$-$7}]}
\newcommand{\agHaikuBN}{259}
\newcommand{\agHaikuBBreakeven}{0.41}
\newcommand{\agHaikuBVtwoEsc}{100\%}
\newcommand{\agHaikuBVtwoErr}{0\%}
\newcommand{\agHaikuBCRej}{15\%}
\newcommand{\agOpusAUEscalate}{94\%}
\newcommand{\agOpusAUAuto}{6\%}
\newcommand{\agOpusAUUnsafeaccept}{1\%}
\newcommand{\agOpusAUWrongreject}{2\%}
\newcommand{\agOpusAUAutoerror}{3\%}
\newcommand{\agOpusAUAutoacc}{0.53}
\newcommand{\agOpusACEscalate}{100\%}
\newcommand{\agOpusACAuto}{0\%}
\newcommand{\agOpusACUnsafeaccept}{0\%}
\newcommand{\agOpusACWrongreject}{0\%}
\newcommand{\agOpusACAutoerror}{0\%}
\newcommand{\agOpusACAutoacc}{--}
\newcommand{\agOpusADEscalate}{\mbox{$+$6.2}}
\newcommand{\agOpusADEscalateCi}{[\mbox{$+$4}, \mbox{$+$9}]}
\newcommand{\agOpusADUnsafeaccept}{\mbox{$-$0.7}}
\newcommand{\agOpusADUnsafeacceptCi}{[\mbox{$-$2}, \mbox{$+$0}]}
\newcommand{\agOpusADAutoerror}{\mbox{$-$2.9}}
\newcommand{\agOpusADAutoerrorCi}{[\mbox{$-$5}, \mbox{$-$1}]}
\newcommand{\agOpusAN}{274}
\newcommand{\agOpusABreakeven}{0.47}
\newcommand{\agOpusACRej}{0\%}
\newcommand{\agOpusAConVsPla}{\mbox{$-$1.8}}
\newcommand{\agOpusAConVsPlaCi}{[\mbox{$-$3.6}, \mbox{$-$0.4}]}
\newcommand{\agOpusAUnpVsNo}{\mbox{$-$1.1}}
\newcommand{\agOpusAUnpVsNoCi}{[\mbox{$-$5.8}, \mbox{$+$3.1}]}
\newcommand{\agOpusAConVsNo}{\mbox{$-$4.0}}
\newcommand{\agOpusAConVsNoCi}{[\mbox{$-$8.2}, \mbox{$-$0.8}]}
\newcommand{\agOpusANotoolEsc}{86\%}
\newcommand{\agOpusANotoolErr}{4.0\%}
\newcommand{\agOpusANotoolRej}{1\%}
\newcommand{\agOpusAPlaceboEsc}{96\%}
\newcommand{\agOpusAPlaceboErr}{1.8\%}
\newcommand{\agOpusAPlaceboRej}{2\%}
\newcommand{\agOpusBUEscalate}{91\%}
\newcommand{\agOpusBUAuto}{9\%}
\newcommand{\agOpusBUUnsafeaccept}{0\%}
\newcommand{\agOpusBUWrongreject}{3\%}
\newcommand{\agOpusBUAutoerror}{4\%}
\newcommand{\agOpusBUAutoacc}{0.58}
\newcommand{\agOpusBCEscalate}{100\%}
\newcommand{\agOpusBCAuto}{0\%}
\newcommand{\agOpusBCUnsafeaccept}{0\%}
\newcommand{\agOpusBCWrongreject}{0\%}
\newcommand{\agOpusBCAutoerror}{0\%}
\newcommand{\agOpusBCAutoacc}{--}
\newcommand{\agOpusBDEscalate}{\mbox{$+$9.3}}
\newcommand{\agOpusBDEscalateCi}{[\mbox{$+$3}, \mbox{$+$16}]}
\newcommand{\agOpusBDUnsafeaccept}{\mbox{$-$0.4}}
\newcommand{\agOpusBDUnsafeacceptCi}{[\mbox{$-$1}, \mbox{$+$0}]}
\newcommand{\agOpusBDAutoerror}{\mbox{$-$3.9}}
\newcommand{\agOpusBDAutoerrorCi}{[\mbox{$-$8}, \mbox{$-$1}]}
\newcommand{\agOpusBN}{259}
\newcommand{\agOpusBBreakeven}{0.42}
\newcommand{\agOpusBCRej}{0\%}
\newcommand{\agPmin}{0.30}
\newcommand{\agPmax}{0.76}
\newcommand{\agCalls}{4{,}264}
\newcommand{\agCallsVtwo}{5{,}191}
\newcommand{\agCost}{11.10}


\begin{document}

\title{Can Agents Trust a Replication Verifier?\\
Epistemic Contracts for Scientific Agent Tools}

\author{\IEEEauthorblockN{Luxin Zhang}
\IEEEauthorblockA{Independent Researcher\\ luxin.zhang92@outlook.com}}

\maketitle

\begin{abstract}
Agentic AI systems can delegate the question ``is this scientific claim trustworthy?'' to a verifier tool, yet such tools are evaluated on accuracy rather than on the guarantees a calling agent relies on. We propose a testable methodology for \emph{epistemic contracts} for agent tools: guarantees about what an output may reflect (temporal validity, isolation from the queried paper's labels), what decisions it licenses (base rate, calibration, selective accuracy), and how it is audited. We apply it to \textsc{SciReliGraph}, an MCP replication verifier. Under a standard mutation-testing tool, its contract tests reach a verified mutation score of \mutVoneScore{}, and \mutVthreeScore{} after tests motivated by surviving mutants. On the 274 SCORE claims with published replication outcomes, a contract-compliant verifier discriminates only modestly (AUROC \aucBone{}), and bibliographic evidence-graph features yield no statistically established improvement. Of four LLM judges, the best rank competitively, but their temporal validity cannot be certified, and a decoy-calibrated probe shows two of the four naming outcomes for papers that were never replicated. In an agent-in-the-loop experiment with no-tool and placebo controls, a bare probability induced near-chance automatic decisions that most agents did not make without the tool; the contract removed most of them, beyond a generic warning, for three of four LLM agents, while the fourth misread it and rejected most claims until the wording stated the intended action. Scientific agents should treat replication forecasts as audited weak signals, and tool contracts should be tested against the agents that consume them.
\end{abstract}

\begin{IEEEkeywords}
agentic AI, epistemic contracts, AI for science, verification, LLM-as-a-judge, contract testing, metamorphic testing, data leakage, MCP, replicability
\end{IEEEkeywords}

\section{Introduction}
Systems that automate parts of the research process---from literature-grounded hypothesis generation to autonomous experimentation~\cite{lu2024}---repeatedly need to decide which prior findings are solid enough to build on. In multi-agent designs such decisions can be delegated to a \emph{critic} or \emph{verifier} tool, for instance an LLM-as-a-judge~\cite{zheng2023}. The decision matters: large coordinated projects found that a substantial fraction of published social and behavioural science findings do not replicate~\cite{osc2015,camerer2016,camerer2018}, and the DARPA SCORE program has released outcomes for 274 independently replicated claims from 164 papers, of which 151 (55.1\%) replicated~\cite{tyner2026,score_g5sny}. Recent evidence also counsels caution about delegating it. LLMs that predict experimental outcomes cannot tell their reliable predictions from their unreliable ones~\cite{sehwag2026scipredict}, and AI-led teams reproduced far fewer published results than human-only or AI-assisted teams~\cite{brodeur2026aiteams}.

An agent that calls a verifier at time $t$ about claim $c$ implicitly assumes that (i)~the answer uses only information available at $t$, (ii)~the returned probability means what it says, and (iii)~the answer can be audited later. We call such guarantees an \emph{epistemic contract}. Security and correctness checks ask whether a tool is safe and runs as specified. An epistemic contract instead states what the tool's \emph{output may reflect} and what decisions it \emph{licenses}. Its ingredients are not new individually: design by contract~\cite{meyer1992}, output documentation~\cite{mitchell2019}, calibration~\cite{guo2017}, metamorphic~\cite{chen2018mt} and mutation testing~\cite{jia2011}, and tool-use evaluation~\cite{wu2026mcpmark} each have long research histories. Our contribution is a \emph{testable methodology} that combines them for agent tools, together with evidence of a failure mode that only becomes visible when they are combined. These guarantees fail silently. A verifier fitted on bibliometric data scraped today can absorb post-replication signals; an LLM judge may simply remember the outcome~\cite{sainz2023}; a random-split evaluation lets claims from one paper appear in both training and test data~\cite{kaufman2012,kapoor2023}. None of these failures is visible to a calling agent.

We ask: \emph{if a verifier is built and tested against an explicit epistemic contract, is it good enough for an agent to rely on, and does the agent use it correctly?} We began by hypothesising that time-bounded evidence graphs (H1), corrected for non-independent support (H2), would improve forecasts, and that lifting the evidence cutoff would inflate apparent accuracy (H3). The data supported none of these, and we report the negative results in full. A further result motivates testing the \emph{consumer} of the contract, not only the tool. In our agent experiment, one LLM agent read the tool's honest warning that its score was weak as evidence \emph{against the claim}, and rejected \agHaikuACRej{} of claims instead of escalating them (\S\ref{sec:agent}).

\textbf{Contributions.}
\begin{enumerate}
\item \emph{A testable epistemic-contract methodology for verifier tools} (\S\ref{sec:contract}): five requirements, a feature-availability typing that makes temporal validity checkable, and a reusable machine-readable specification (a JSON Schema for answers and refusals, with tool-agnostic test relations). In our own pipeline, the typing caught a leak that would otherwise have produced our ``best'' harness result.
\item \emph{Contract tests with verified adequacy} (\S\ref{sec:harness}--\S\ref{sec:tool}): an MCP tool that enforces the contract at call time, metamorphic and isolation tests, and a two-level mutation analysis (targeted mutants, and a standard tool whose verdicts we re-verified independently).
\item \emph{A decoy-calibrated memorization probe for LLM judges} (\S\ref{sec:llm}). Asking judges about journal-matched papers that were never replicated or reproduced separates recall from fabrication: two of four models named outcomes for such papers.
\item \emph{Evidence of a contract--consumer gap} (\S\ref{sec:agent}). With no-tool and placebo controls, a bare verifier probability induced errors that most agents did not make without the tool; the contract removed most of them, beyond a generic warning, for three agents, and misled a fourth, which in an exploratory follow-up was fixed by stating the intended action.
\end{enumerate}
Empirically, on the full SCORE replication set, the best leakage-safe verifiers are weak (AUROC $\approx$0.63), bibliographic evidence graphs yield no statistically established improvement, and the best LLM judges are competitive but uncertifiable. Replication forecasts should therefore inform agents, not gate them.

\section{Background and Related Work}
\textbf{Replication outcomes and their prediction.} Large replication projects~\cite{osc2015,camerer2016,camerer2018} were accompanied by surveys and prediction markets. In a pooled analysis of 103 findings, prediction markets classified replication outcomes with 73\% accuracy~\cite{gordon2021,dreber2015}. For SCORE, two large human-forecasting efforts reached a median AUC of 0.76, validated on a subset of 37 high-powered replications~\cite{mody2026}. Machine-learning forecasters using statistical and textual features of original papers have also reported above-chance accuracy~\cite{altmejd2019,yang2020,youyou2023}. SCORE additionally produced machine credibility assessments, including feature-based and knowledge-graph-reasoning systems, developed largely to predict \emph{human} assessments, with a limited set of replication outcomes for validation~\cite{alipourfard2021,rajtmajer2026}. Hybrid prediction markets in which algorithmic agents trade alongside humans have also been studied~\cite{chakravorti2026hybrid}, and ReplicatorBench evaluates LLM agents that \emph{carry out} replications rather than forecast them~\cite{nguyen2026replicatorbench}. Our work differs in three respects. We use only realised replication outcomes as labels. We treat temporal leakage as an enforced requirement rather than a convention. And we evaluate the verifier as a tool whose output an agent will act on.

\textbf{Evidence structure.} Nonreplicable findings have been reported to be cited more~\cite{serragarcia2021}; we test whether citation \emph{structure} helps where raw counts do not. Unlike fact-checking~\cite{wadden2020}, which asks whether existing documents support a claim, replication forecasting predicts a future study and is exposed to hindsight leakage.

\textbf{Calibration and verification of LLM judgements.} Evaluations of LLM forecasters are vulnerable to many forms of temporal leakage~\cite{paleka2026pitfalls}, which motivates our enforced cutoffs and memorization probe. GLEAN calibrates trajectory evidence into correctness probabilities for high-stakes agent verification and triggers extra verification under uncertainty~\cite{zhang2026glean}. Claim-level calibration applies the same post-hoc calibration to every baseline before comparing them~\cite{abbasli2026claimcal}. We follow that practice when recalibrating our LLM judges.

\textbf{Verifying MCP tools and agents.} MCP-specific work evaluates agents' multi-step tool use~\cite{wu2026mcpmark}, explores MCP tool spaces with metamorphic evaluation~\cite{liu2026sgvef}, defends the tool layer against poisoning and prompt injection at run time~\cite{yergattikar2026securing}, and records agents' MCP actions to evaluate experiment reproduction as a process~\cite{hong2026agentactionbench}. These works ask whether an agent \emph{uses} tools correctly or safely. We ask a complementary question: whether a tool's \emph{output} keeps an epistemic contract that the calling agent can rely on.

\begin{table}[t]\centering
\caption{Positioning against the closest work, by each work's stated scope. \checkmark: a stated focus; $\circ$: addressed in a different sense; --: not a stated focus (not a claim of absence). Our agent evaluation is single-turn, hence $\circ$ in (e). Columns: (a)~realised outcomes as labels; (b)~temporal leakage controlled; (c)~calibration or selective accuracy evaluated; (d)~the tool's output contract tested; (e)~agent use of the output evaluated.}\label{tab:position}
\scriptsize\setlength{\tabcolsep}{3pt}\begin{tabular}{@{}lccccc@{}}\toprule
Work & (a) & (b) & (c) & (d) & (e) \\\midrule
Replication predictors~\cite{altmejd2019,yang2020,youyou2023,rajtmajer2026} & \checkmark & -- & -- & -- & -- \\
SciPredict~\cite{sehwag2026scipredict} & \checkmark & \checkmark & \checkmark & -- & -- \\
GLEAN~\cite{zhang2026glean}, CLAIM-CAL~\cite{abbasli2026claimcal} & \checkmark & -- & \checkmark & -- & $\circ$ \\
ReplicatorBench~\cite{nguyen2026replicatorbench} & \checkmark & -- & -- & -- & \checkmark \\
MCP agent benchmarks~\cite{wu2026mcpmark,liu2026sgvef,hong2026agentactionbench} & -- & -- & -- & $\circ$ & \checkmark \\
MCP runtime defence~\cite{yergattikar2026securing} & -- & -- & -- & $\circ$ & -- \\
\textbf{This work} & \checkmark & \checkmark & \checkmark & \checkmark & $\circ$ \\
\bottomrule\end{tabular}\end{table}

Table~\ref{tab:position} summarises the gap: prior work evaluates verifiers' accuracy or agents' tool use, but not whether a verifier's output keeps a testable contract \emph{and} whether agents consume that contract correctly.

\textbf{SE for ML and agent tools.} Hidden data dependencies are a recognised source of ML technical debt~\cite{sculley2015,amershi2019}, and datasheets~\cite{gebru2021} make provenance explicit. MCP~\cite{mcp2024} makes it easy to expose a predictor as a tool, but not to specify or test what its output may be trusted for. We contribute a contract, contract tests, and evidence of their adequacy for one safety-relevant tool class.

\section{A Verifier Contract for Agentic Pipelines}\label{sec:contract}
\subsection{Problem definition}
Let $c$ be an original empirical claim, $t$ a forecast cutoff, and $\mathcal{E}(c,t)$ the admissible evidence dated $\le t$. Let $Y_c\in\{0,1\}$ be the outcome of a later independent replication under a pre-declared criterion. The verifier returns $\hat p_c = \hat P(Y_c=1\mid c,\mathcal{E}(c,t))$ with an evidence record. We evaluate discrimination (AUROC), probabilistic accuracy with strictly proper scores (Brier, log loss)~\cite{gneiting2007}, calibration~\cite{guo2017}, and selective accuracy~\cite{elyaniv2010}. The last of these matters because an agent may act only on confident outputs.

\subsection{Requirements of the epistemic contract}
\textbf{R1 Temporal validity.} No feature or evidence of the queried claim may encode information dated after $t$, in particular replication results, replication-paper metadata, or post-cutoff citations. (The model is necessarily trained on \emph{other} claims' later outcomes, and the reported base rate is itself post-$t$; we treat this residual exposure as a limitation.) Calls whose $t$ falls after the outcome may be public are refused.
\textbf{R2 Label validity.} The label is the declared criterion, reconciled to published counts; alternative criteria are reported as sensitivity analyses.
\textbf{R3 Honest generalisation.} Claims from the same paper are never split across training and test data, and a claim is never scored by a model that saw any label from its own paper.
\textbf{R4 Decision-relevant output.} The output must beat the base rate on proper scores and must expose the base rate. If it is used for gating, its confident outputs must be more accurate than its unconfident ones.
\textbf{R5 Auditability.} Every prediction must be traceable to a source archive (with checksum), an evidence access date, the feature tiers used, and a model configuration.
R1 and R3 constrain what the output may \emph{reflect}, R4 what it \emph{licenses}, and R2 and R5 make both checkable. R1, R2, R3 and R5 are enforceable integrity requirements. R4 is different: it is an empirical acceptance criterion that a verifier may fail. A contract-compliant response must then \emph{expose} that failure, through its base rate, skill and gating verdict, rather than imply that the predictor is decision-ready. In our results, R1--R3 and R5 are met by construction and tested (\S\ref{sec:tool}). R4's first clause is met only nominally (B1 Brier skill \bssBone{}, paired CI includes zero), and its gating clause is not met by any harness verifier (\S\ref{sec:selective}).

\subsection{Feature-availability typing}
R1 is enforced by \emph{typing} every feature with the moment it becomes knowable.
\textbf{Tier A} (original publication): sample size; $p$-value and whether it is reported exactly; marginal significance ($0.01<p\le0.05$); converted effect size $|r|$; power to detect half the original effect; interaction vs.\ main effect; test-statistic and analysis family; publication year; and discipline.
\textbf{Tier T} (original publication): title, abstract and keywords.
\textbf{Tier G($t$)} (cutoff-dependent): the evidence graph of \S\ref{sec:harness}.
\textbf{Tier D} (replication launch): new vs.\ secondary data, the preregistered power target, and the \emph{planned} sample size implied by that target.
\textbf{Tier D$^{*}$} (after data collection): the \emph{achieved} replication sample size and the power computed from it.
A forecast ``at publication'' may use A, T and G($t$); a forecast ``at launch'' may add D. D$^{*}$ is never admissible and appears only as a leakage illustration: the achieved size exceeds the planned size by a median factor of \plannedRatio{} and can depend on how data collection went. Typing also forced decisions about fields whose timing the archive does not record. SCORE's journal-level TOP Factor policy scores look like publication metadata, but their rating date is not stored, so R1 cannot be verified for them and we exclude them.

\section{The \textsc{SciReliGraph} Harness}\label{sec:harness}
% Pipeline diagram removed for space; Table req maps each requirement to its enforcement.

\begin{table}[t]\centering
\caption{Requirements, enforcement mechanisms, and what enforcement caught or showed in this study.}\label{tab:req}
\scriptsize\setlength{\tabcolsep}{2.5pt}
\begin{tabular}{@{}p{0.11\columnwidth}p{0.42\columnwidth}p{0.41\columnwidth}@{}}\toprule
Req. & Enforcement mechanism & Finding in this study \\\midrule
R1 temporal & Feature tiers; dated-evidence filter; assertion that the cutoff precedes the earliest replication start; call-time refusal; metamorphic test & Excluded TOP Factor scores (undated); exposed the achieved-size leak (Tier D$^{*}$) \\
R2 label & Build-time assertion of 274/164/151; alternative criteria kept as separate labels & Exact reconciliation; three sensitivity labels \\
R3 generalisation & Paper-grouped CV; nested grouped tuning; leave-paper-out in the tool; label-flip test & Unregularised starter baseline scored worse than the base rate \\
R4 decision use & Proper scores vs.\ base rate; reliability diagram (artifact); selective accuracy; \texttt{selective\_ok} flag & B1 Brier skill \bssBone{} (CI includes 0); no harness verifier shows a supported gating gain \\
R5 audit & Manifest (sha256, access date, tiers, missingness); per-call audit log; scripted tables and numbers & All reported numbers regenerate from outputs \\
\bottomrule\end{tabular}\end{table}

Table~\ref{tab:req} maps each requirement to its enforcement point: the label gate (R2), the evidence cutoff (R1) and the tool boundary (R1, R3--R5). The implementation is \locPython{} lines of Python and runs end to end from public sources.

\textbf{Data and E0 gate.} The harness downloads the publication-frozen SCORE analyst package~\cite{score_g5sny} and records its SHA-256. It selects version-of-record, non-COVID replications and \emph{asserts} that they reproduce the published totals of 274 claims, 164 papers and 151 successes~\cite{tyner2026}. Of the plausible selections we tried (e.g., 427 reports, 309 claims including COVID-era studies), only this one reproduces the publication. The set spans 54 journals, publication years 2009--2018 and six fields; 128 replications used new data and 146 secondary data.

\textbf{Evidence and cutoffs.} For each original paper, resolved by DOI (164/164), the harness caches the OpenAlex~\cite{priem2022} record and every citing work: \nCiting{} in total, with authorships, venues, topics and references. Time filtering is applied only at feature construction. The \emph{strict} cutoff, 2018-12-31, is asserted to precede the earliest replication start in the archive (\earliestStart{}); \nCitingStrict{} citing works fall before it. The \emph{launch} cutoff is the day before each replication project started; it is known for \nLaunch{} claims, and the remainder use the strict cutoff. The \emph{none} cutoff admits everything available at access time (October 2026) and is a diagnostic only.

\textbf{Evidence graph G($t$).} For claim $c$, the graph contains the original paper and all citing works dated $\le t$. It has three edge types: \emph{cites} (citer$\rightarrow$original and citer$\rightarrow$citer), \emph{shares-author} (a dependency edge) and \emph{shares-reference} (bibliographic coupling). From it we derive three feature groups.
(i)~\emph{Flat evidence}: log citations, citations per year since publication, and log reference count.
(ii)~\emph{Graph structure}: citing-venue entropy, the number of citing fields, the cross-field citer share, citer--citer citation density, mean citer--original coupling, and an evidence indicator.
(iii)~\emph{Dependency correction}: self-citation share, the number of independent support groups (connected components of non-self citers under shared authorship), the ratio of groups to citers, and the number of distinct citing authors.
Two caveats apply. The graph is \emph{bibliographic}: we did not extract semantic supports/contradicts edges. And its features are paper-level, so all claims from one paper share them.

\textbf{Models and protocol.} B0 predicts the training-fold prevalence. Tabular models use median imputation with missingness indicators, standardisation and one-hot encoding (rare levels pooled), followed by L2 logistic regression. The regularisation strength $C\in\{0.003,\dots,1\}$ is chosen by inner grouped 3-fold CV on log loss. A gradient-boosted tree baseline (depth 2, 150 iterations) has fixed hyperparameters. The text model is TF-IDF (1--2-grams) with tuned logistic regression. We use stratified, paper-grouped 5-fold CV repeated over 20 seeds, and report the mean and SD over repetitions. For inference we average each claim's out-of-fold predictions over seeds and draw 2{,}000 paper-clustered bootstrap resamples. Paired differences are computed on identical resamples and Holm-adjusted over the family of 11 planned comparisons. The \textbf{permutation null} permutes all graph and dependency features across papers within publication-year strata, \permK{} times, re-running the full CV each time. This preserves marginal distributions and the year--citation relation while breaking the link to outcomes.

\section{The Verifier as an Agent Tool}\label{sec:tool}
Offline evaluation establishes how good a verifier is; it does not ensure that the deployed tool behaves as evaluated. We therefore implemented the verifier as a library with an MCP~\cite{mcp2024} server exposing a single tool, \texttt{verify\_claim(claim\_id, as\_of, stage)}, and enforced the contract inside the call path.
(1)~\emph{R1:} for every call, the graph features of \emph{all} claims are recomputed from the raw cache with evidence filtered to \texttt{as\_of}, and calls on or after the claim's replication launch are refused.
(2)~\emph{R3:} the model is refitted without any claim from the queried paper (leave-paper-out), with the same nested tuning as the harness.
(3)~\emph{Stage gating:} \texttt{publication} uses tiers A+G, and \texttt{launch} uses A+D.
(4)~\emph{R4:} the response carries the base rate, the cross-validated AUROC and Brier skill of that configuration, and a \texttt{selective\_ok} flag that is true only if the lower 95\% bound of top-decile accuracy exceeds overall accuracy in the selective-accuracy test (\S\ref{sec:selective}); it is false for every harness configuration we evaluated.
(5)~\emph{R5:} each call appends a record with inputs, outputs, archive hash and model-configuration hash to an audit log.
Listing~\ref{lst:tool} shows the actual responses an MCP client received in our smoke test: one valid call and one refused call.

\lstinputlisting[caption={Actual \texttt{verify\_claim} responses over MCP stdio (truncated hashes).},label={lst:tool},float=t]{tables/tool_output.txt}

\textbf{Contract tests and their adequacy.} Unit tests that check outputs against expected values do not apply: the ``right'' probability is unknown. We therefore test \emph{relations} that must hold under the contract (Table~\ref{tab:contract}).
\begin{itemize}
\item A \emph{metamorphic} test injects synthetic citing works dated after \texttt{as\_of} into every paper's evidence and requires a bit-identical output. A paired positive control checks that works dated before \texttt{as\_of} \emph{are} counted.
\item An \emph{isolation} test flips the labels of every claim from the queried paper and requires an unchanged prediction.
\item A \emph{gating} test randomises Tier-D values and requires an unchanged publication-stage output.
\item Further tests check refusals, malformed calls, the reported base rate, and audit-record completeness.
\end{itemize}
To check that these tests would actually detect contract violations, we built \nMutants{} mutants. Each removes one guarantee: evidence filtering, leave-paper-out, stage gating, the launch refusal, or audit logging. The unmodified tool passes all \nTests{} original tests, and every mutant is killed by at least one test (\nKilled{}/\nMutants{}). The positive control kills no mutant by design; its role is to make the metamorphic test non-vacuous. The response shape is fixed by a JSON Schema that any verifier can adopt; schema-conformance tests check the tool's answers and refusals against it.

\textbf{Independent mutation analysis.} Because we wrote these five mutants ourselves, we also ran a standard tool, \texttt{mutmut}~2.5.1, over the verifier module. It generated \mutTotal{} mutants (Table~\ref{tab:mutstages}). Its own verdicts were unreliable in our environment, missing some kills. We therefore re-ran every reported survivor in isolation, with a freshly restored source file and no bytecode cache, and \mutMissed{} turned out to be killed. The remaining survivors exposed real gaps, closed in two rounds:
\begin{itemize}
\item \emph{Round 1, ten tests:} the exact output schema, the true branch of the gating rule, evidence and training-set counts, and audit-record hashing.
\item \emph{Round 2, four tests:} the temporal boundary that R1 depends on and two oracle checks. A work dated exactly $t$ must be admitted and one dated $t{+}1$ excluded; claims without a launch date must fall back to the archive's earliest replication start; and loaded evidence must match the raw cache files.
\end{itemize}
Because these fourteen tests were written after seeing the survivors, the final score is an optimistic measure of adequacy, and we report it alongside the pre-analysis \mutVoneScore{}. None of the \mutSurv{} remaining survivors threatens a contract guarantee: \mutClsModelHyperparameters{} change model hyperparameters that the contract deliberately leaves open, \mutClsCliEntryPoint{} touch the command-line entry point, \mutClsExplanationDrivers{} alter the explanatory ``drivers'' output, which no test checks, \mutClsMessageOrRounding{} change message wording or rounding, and \mutClsEquivalent{} are equivalent or near-equivalent.



\begin{table*}[t]\centering\caption{Contract tests for the verifier tool and mutation analysis. Each mutant removes one guarantee from the implementation; $\times$ marks the tests that fail (kill the mutant). The unmodified verifier passes all 9 tests.}\label{tab:contract}
\footnotesize\setlength{\tabcolsep}{4pt}\begin{tabular}{@{}lp{0.52\textwidth}ccccc@{}}\toprule
Req. & Contract test & leaky & no-LPO & no-gate & no-refuse & no-audit \\\midrule
R1 & Metamorphic: citing works dated after \texttt{as\_of} added to every paper leave output unchanged & $\times$ &  &  &  &  \\
R1 & Positive control: works dated before \texttt{as\_of} are counted &  &  &  &  &  \\
R1 & Call on/after the replication launch is refused &  &  &  & $\times$ &  \\
-- & Malformed date or unknown stage is refused (2 cases) &  &  &  &  &  \\
R3 & Flipping all labels of the queried claim's paper leaves its prediction unchanged &  & $\times$ &  &  &  \\
R1 & Publication stage: Tier-D values have no influence &  &  & $\times$ &  &  \\
R4 & Base rate equals training prevalence; weak-signal advice when selective test fails &  & $\times$ &  &  &  \\
R5 & Each call appends an audit record with provenance hashes &  &  &  &  & $\times$ \\
\midrule
\multicolumn{2}{@{}l}{Mutant killed?} & yes & yes & yes & yes & yes \\
\bottomrule\end{tabular}\end{table*}

\begin{table}[t]\centering
\caption{Mutation score of the verifier module (\textnormal{\texttt{mutmut}}, \mutTotal{} mutants), by test suite. Each verified row re-runs all surviving mutants in isolation. The three later schema-conformance tests were not part of these runs.}\label{tab:mutstages}
\scriptsize\setlength{\tabcolsep}{4pt}\begin{tabular}{@{}lccc@{}}\toprule
Stage & Tests & Killed & Score \\\midrule
Tool's own report & 9 & \mutToolKilled{} & \mutToolScore{} \\
Verified: original contract tests & 9 & \mutVoneKilled{} & \mutVoneScore{} \\
\quad + round 1 (schema, gating rule, counts, audit) & 19 & \mutVtwoKilled{} & \mutVtwoScore{} \\
\quad + round 2 (cutoff boundary, oracles) & 23 & \mutVthreeKilled{} & \mutVthreeScore{} \\
\quad excluding \mutClsEquivalent{} equivalent mutants & 23 & \mutVthreeKilled{} & \mutVthreeExEq{} \\
\bottomrule\end{tabular}\end{table}

\section{Results on the SCORE Testbed}\label{sec:results}
Replication forecasting serves here as the testbed for the contract: it supplies realised outcomes, a natural temporal boundary, and a decision an agent would plausibly delegate. We first ask how good a contract-compliant verifier is (A--E), then how LLM judges behave under the same contract (F), and finally check robustness (G).
\begin{table*}[t]\centering\caption{Out-of-fold forecasting performance on 274 SCORE claims (164 papers). AUROC: mean$\pm$SD over 20 repetitions of paper-grouped 5-fold CV, with paper-clustered bootstrap 95\% CI of the seed-averaged predictions. Brier: mean over repetitions (lower is better). BSS: Brier skill score vs.\ B0. Cutoff ``post'': uses information available only after replication data collection (leakage illustration).}\label{tab:main}
\footnotesize\setlength{\tabcolsep}{6pt}\begin{tabular}{@{}llccccc@{}}\toprule
ID & Features & Cutoff & AUROC & 95\% CI & Brier & BSS \\\midrule
B0 & Training-fold prevalence & strict & 0.492$\pm$0.001 & [0.37, 0.56] & 0.2474 & +0.0\% \\
B1 & Original-study metadata, LR & strict & 0.627$\pm$0.017 & [0.55, 0.71] & 0.2409 & +2.6\% \\
B1g & Original-study metadata, GBT & strict & 0.620$\pm$0.020 & [0.55, 0.72] & 0.2483 & -0.3\% \\
B2 & Title/abstract TF-IDF, LR & strict & 0.481$\pm$0.022 & [0.37, 0.55] & 0.2487 & -0.5\% \\
B4 & B1 + flat citation counts & strict & 0.613$\pm$0.019 & [0.54, 0.69] & 0.2432 & +1.7\% \\
M1 & B4 + graph structure & strict & 0.624$\pm$0.018 & [0.55, 0.70] & 0.2406 & +2.8\% \\
M2 & M1 + dependency correction & strict & 0.620$\pm$0.020 & [0.55, 0.70] & 0.2406 & +2.8\% \\
G & Graph + dependency only & strict & 0.591$\pm$0.027 & [0.53, 0.69] & 0.2485 & -0.4\% \\
\midrule
NC & M2, graph permuted (1 draw) & strict & 0.611$\pm$0.019 & [0.53, 0.69] & 0.2421 & +2.2\% \\
\midrule
M2$_\ell$ & M2 & launch & 0.620$\pm$0.020 & [0.55, 0.70] & 0.2405 & +2.8\% \\
D & B1 + planned replication design & launch & 0.625$\pm$0.019 & [0.54, 0.71] & 0.2429 & +1.8\% \\
D$^{*}$ & D + achieved size/power & post & 0.638$\pm$0.024 & [0.58, 0.72] & 0.2403 & +2.9\% \\
\midrule
M2$_\infty$ & M2 & none & 0.614$\pm$0.017 & [0.53, 0.71] & 0.2409 & +2.7\% \\
G$_\infty$ & Graph + dependency only & none & 0.579$\pm$0.020 & [0.49, 0.68] & 0.2466 & +0.3\% \\
\bottomrule\end{tabular}\end{table*}

% Paired-differences table (paper/tables/delta.tex) moved to the artifact for space; all comparisons are quoted in the text.

\subsection{How good is a careful verifier?}
Table~\ref{tab:main} summarises all configurations. The publication-time metadata model B1 reaches AUROC \aucBone{}$\pm$\sdBone{}, and its CI \ciBone{} excludes 0.5. B0's AUROC below 0.5 is an artefact of grouped CV (removing a fold's papers shifts the training prevalence opposite to that fold's own), so comparing AUROCs with B0 is not meaningful; the relevant evidence of signal is B1's CI excluding 0.5. The probabilistic gain is small: Brier \brierBone{} vs.\ \brierBzero{}, a skill score of \bssBone{}, with a paired CI that includes zero (paired bootstrap; full table in the artifact). The gradient-boosted variant ranks similarly but is worse calibrated.

The individual signals agree with the replication literature. Replicated claims had stronger original evidence (mean $-\log_{10}p$ \fpRep{} vs.\ \fpNon{}; single-feature AUROC \fpAuc{}), higher original power (\fpowRep{} vs.\ \fpowNon{}), and less often fell in the $0.01<p\le0.05$ band (\fmargRep{} vs.\ \fmargNon{}). Even combined, these signals give only modest discrimination.

An unregularised starter baseline with one-hot journals scored \emph{worse} than the base rate (Brier \starterBrier{} vs.\ \starterPrevBrier{}), so base-rate comparison belongs in a verifier's acceptance tests.

% AUROC figure (paper/figures/auroc.pdf) omitted for space; Table main carries the same CIs.

% Full single-feature table: paper/tables/feat.tex (released with the artifact)

\subsection{H1/H2: Do evidence graphs and dependency correction help?}
Not detectably. Flat citation evidence added to metadata changes AUROC by \dBfourBone{} (CI \dciBfourBone{}; $p_{\mathrm{Holm}}=\holmBfourBone{}$). Graph structure on top of flat evidence (M1 vs.\ B4: \dMoneBfour{}, CI \dciMoneBfour{}) and dependency correction on top of that (M2 vs.\ M1: \dMtwoMone{}, CI \dciMtwoMone{}) are both null. Overall, M2 is indistinguishable from metadata alone (\dMtwoBone{}, CI \dciMtwoBone{}).

The permutation null sharpens this picture. With the same \permSeeds{} CV seeds, M2 reaches AUROC \permReal{}. Over \permK{} within-year permutations of the graph block, the null mean is \permNullMean{} (2.5--97.5\% range \permNullLo{}--\permNullHi{}), so the real graph features sit at the edge of the null distribution ($p=\permP{}$). They may carry weak signal, but not enough to offset the variance they add: on the same seeds B1 reaches \permBone{}. Alone, graph features are weak (AUROC \aucGonly{}). The strongest single descriptor is the share of \emph{cross-field} citers, which is lower for replicated claims (\fxfRep{} vs.\ \fxfNon{}; single-feature AUROC \fxfAuc{} in the reverse direction). In the multivariable models this adds nothing to metadata, possibly because it reflects field composition. Raw pre-cutoff citation counts are essentially uninformative (AUROC \fcitesAuc{}), in line with~\cite{serragarcia2021}.

\subsection{H3: Does removing the cutoff inflate performance?}
Not detectably for these features. Admitting all evidence available in 2026 raises the mean $\log(1+\text{citations})$ from \logCitesStrict{} to \logCitesLeaky{}, and evidence coverage from \hasEvStrict{} to \hasEvLeaky{}. Discrimination, however, does not improve (M2$_\infty$ vs.\ M2: \dLeakMtwo{}, CI \dciLeakMtwo{}; G$_\infty$ vs.\ G: \dLeakG{}, CI \dciLeakG{}). Because these features carry little signal even when leaky, this test has low power. It does not license dropping the cutoff, and it says nothing about text or LLM verifiers, whose training data may contain the outcomes~\cite{sainz2023}.

\subsection{Launch-time information and a leak we caught}
Moving the cutoff to each replication's start date leaves performance unchanged (M2$_\ell$: \aucMtwolaunch{}). Our first ``launch-time'' model used the archive's replication sample size and power and appeared best; typing revealed these are \emph{achieved} values, known only after data collection. With the planned design only (Tier D), the gain disappears (D vs.\ B1: \dDplanBone{}, CI \dciDplanBone{}). Adding the achieved values back (D$^{*}$) recovers part of it (\dDachDplan{}, CI \dciDachDplan{}). The difference is not significant, but its direction is what post-hoc information would produce. The lesson for tool builders is that a field's \emph{name} (``replication sample size'') does not reveal its availability time; the typing step did.

% Harness selective-accuracy table (paper/tables/selective.tex) reported in text for space.

\subsection{Can an agent gate on confidence?}\label{sec:selective}
An agent may act only when a verifier is confident. We rank claims by confidence $|p-0.5|$ and report accuracy on the most confident fraction (full table in the artifact); for a verifier suitable for gating, accuracy should rise as coverage falls. The most confident 10\% of B1's predictions are correct \selTopBone{} of the time (CI \selTopCIBone{}), versus \selAllBone{} over all claims; for D the figures are \selTopDplan{} vs.\ \selAllDplan{}. M2's top decile is more accurate than its full output (\selTopMtwo{} vs.\ \selAllMtwo{}), but its interval \selTopCIMtwo{} is far too wide to support gating. Between \selPosMin{} and \selPosMax{} of these top-decile predictions are confident \emph{replicate} calls. At 50\% coverage accuracy is slightly higher than overall (B1 \selHalfBone{}, M2 \selHalfMtwo{}), but the CIs overlap broadly. The reliability diagram (in the artifact) explains why: mean predictions per quintile span a narrow range, and for B1 the top quintile replicates less often than the fourth. R4's gating clause is therefore not met by any harness configuration.

% Reliability diagram (paper/figures/calibration.pdf) omitted for space.

\begin{table*}[t]\centering\caption{LLM-as-judge verifiers on the same claims. \emph{Identified}: title, citation, abstract and original statistics; \emph{de-identified}: field and original statistics only. AUROC with paper-clustered 95\% CI; $\Delta$ vs.\ B1 on identical claims and resamples. Probe: share answering that the paper was replicated in SCORE (\emph{aware}) for real SCORE papers vs.\ journal-matched decoys that were never replicated, and accuracy of named outcomes (counted whenever an outcome is named, regardless of stated awareness) vs.\ the majority-outcome baseline on the same claims. Brier (Platt): after Platt scaling fitted on training folds only (paper-grouped CV, 20 repetitions). B1 here is scored on its seed-averaged predictions, so its Brier differs slightly from the per-seed mean in Table~\ref{tab:main}.}\label{tab:llm}
\footnotesize\setlength{\tabcolsep}{4pt}\begin{tabular}{@{}llccccccc@{}}\toprule
Model & Input & AUROC & 95\% CI & $\Delta$AUROC vs.\ B1 & 95\% CI & Brier & Brier (Platt) & Top-10\% acc. \\\midrule
Claude Haiku 4.5 (2025-07) & identified & 0.566 & [0.48, 0.65] & -0.061 & [-0.154, +0.026] & 0.2541 & 0.2466 & 0.52 \\
 & de-identified & 0.592 & [0.51, 0.68] & -0.036 & [-0.132, +0.057] & 0.2561 & 0.2479 & 0.70 \\
Claude Opus 5.5 (2026-06) & identified & 0.655 & [0.57, 0.74] & +0.027 & [-0.044, +0.103] & 0.2392 & 0.2323 & 0.89 \\
 & de-identified & 0.620 & [0.53, 0.70] & -0.007 & [-0.097, +0.077] & 0.2478 & 0.2393 & 0.70 \\
GPT-5.4 mini (2025-08) & identified & 0.600 & [0.52, 0.68] & -0.028 & [-0.106, +0.050] & 0.2753 & 0.2438 & 0.70 \\
 & de-identified & 0.570 & [0.48, 0.66] & -0.058 & [-0.160, +0.038] & 0.2758 & 0.2473 & 0.56 \\
GPT-6.1 sol (2026-04) & identified & 0.663 & [0.58, 0.74] & +0.036 & [-0.031, +0.105] & 0.2594 & 0.2313 & 0.78 \\
 & de-identified & 0.612 & [0.53, 0.70] & -0.016 & [-0.107, +0.073] & 0.2627 & 0.2413 & 0.74 \\
\midrule
B1 (metadata LR) & same claims & 0.628 & [0.55, 0.71] & -- & -- & 0.2387 & -- & 0.44 \\
\midrule
\multicolumn{9}{@{}l}{\emph{Memorization probe}} \\
Model & & Aware (real) & Aware (decoy) & Named outcomes & Named acc. & Majority & & \\\midrule
Claude Haiku 4.5 (2025-07) & & 0.0\% & 0.0\% & 0/274 & -- & -- & & \\
Claude Opus 5.5 (2026-06) & & 0.0\% & 0.0\% & 3/274 & 0.67 & 0.67 & & \\
GPT-5.4 mini (2025-08) & & 42.7\% & 24.0\% & 77/274 & 0.64 & 0.56 & & \\
GPT-6.1 sol (2026-04) & & 2.2\% & 2.0\% & 6/274 & 1.00 & 0.83 & & \\
\bottomrule\end{tabular}\end{table*}


\subsection{LLM-as-judge verifiers}\label{sec:llm}
Agents can also delegate this decision to an LLM, so we evaluated four LLMs from two vendors under the same scoring protocol. Each vendor contributed one model whose training data end before the SCORE results paper (April 2026~\cite{tyner2026}) and one whose data end after it, using vendor-stated cutoffs (the training-data cutoff where the vendor distinguishes it): \texttt{gpt-5.4-mini-2026-03-17} (August 2025) and \texttt{gpt-6.1-sol} (April 2026), \texttt{claude-haiku-4-5-20251001} (training data to July 2025) and \texttt{claude-opus-5-5} (June 2026). Each model received a system prompt dated at the strict cutoff and asking for the probability that a high-powered replication would succeed, as JSON. We used one sample per claim and each model's lowest available reasoning setting. There were three conditions:
(i)~\emph{identified}: title, citation, abstract and the Tier-A statistics;
(ii)~\emph{de-identified}: field and statistics only, which removes the cues needed to recognise the paper;
(iii)~a \emph{memorization probe} asking directly whether the paper was replicated in SCORE and with what outcome, also asked about 100 \emph{decoys}: journal-matched papers that SCORE sampled but never replicated or reproduced.
All responses parsed; no request was refused. The judging calls cost \$\llmCostTotal{} in total.

\emph{Discrimination and calibration} (Table~\ref{tab:llm}). The two post-cutoff models achieved the highest observed AUROCs of all verifiers in this study (GPT-6.1 sol \llmSolIdAuc{}, CI \llmSolIdCi{}; Claude Opus 5.5 \llmOpusIdAuc{}, CI \llmOpusIdCi{}), without a statistically established advantage over B1 (\llmSolIdVsBone{} and \llmOpusIdVsBone{}). The pre-cutoff models were below B1 (\llmMiniIdAuc{} and \llmHaikuIdAuc{}). Raw probabilities differ sharply in quality. Opus 5.5's raw Brier score is as low as B1's (\llmOpusIdBrier{} vs.\ \llmBoneBrier{}; Brier mixes ranking and calibration, so this does not establish equal calibration), whereas the other three are worse than the base rate (\llmSolIdBrier{}, \llmMiniIdBrier{}, \llmHaikuIdBrier{}; B0 \llmBzeroBrier{}). Following~\cite{abbasli2026claimcal}, we applied the same in-fold Platt scaling to every judge. This brings GPT-6.1 sol to \rcSolPlattBrier{} and Opus 5.5 to \rcOpusPlattBrier{} (Brier skill \rcSolPlattBss{} and \rcOpusPlattBss{}), the lowest of any verifier, while leaving AUROC almost unchanged. Neither recalibrated judge is distinguishable from B1 ($\Delta$Brier CI \rcSolPlattVsBoneCi{} and \rcOpusPlattVsBoneCi{}). The pre-cutoff judges improve only marginally (\rcHaikuPlattBrier{} and \rcMiniPlattBrier{} vs.\ B0 \llmBzeroBrier{}). For the post-cutoff judges, the probabilistic failure is mainly one of calibration rather than ranking. Recalibration, however, requires labelled outcomes, so a recalibrated judge is no longer zero-shot.

\emph{Memorization and fabrication.} Removing identifying information lowers the point estimate of AUROC for both post-cutoff models (GPT-6.1 sol \llmSolIdDeid{}, CI \llmSolIdDeidCi{}; Opus 5.5 \llmOpusIdDeid{}, CI \llmOpusIdDeidCi{}; these LLM comparisons were not in the planned Holm family). The pattern is suggestive, not established: the pre-cutoff GPT-5.4 mini shows a similar drop (\llmMiniIdDeid{}), and de-identification also removes the abstract, which is legitimate information. Both rarely recall outcomes explicitly: they name outcomes for only \llmSolNamed{} and \llmOpusNamed{} of 274 claims (GPT-6.1 sol's \llmSolNamed{} are all correct). Their gain from identification therefore reflects either implicit memorization or a legitimate reading of the abstract, which cannot be separated from outside the model; R1 cannot be certified for these judges~\cite{paleka2026pitfalls}. The probe also exposes a model-specific failure. GPT-5.4 mini claims knowledge of a replication for \llmMiniAwareReal{} (\llmMiniAwareRealCi{}) of real papers and \llmMiniAwareDecoy{} (\llmMiniAwareDecoyCi{}) of decoys, and names outcomes for \llmMiniNamedDecoy{} decoys that were never replicated: fabricated replication records. GPT-6.1 sol does so occasionally (\llmSolNamedDecoy{} decoys); both Claude models named none. Unlike instance-level contamination tests~\cite{golchin2024}, the decoys calibrate a model's claimed recall against papers where any claimed outcome must be false. GPT-5.4 mini's excess awareness for real over decoy papers, with named-outcome accuracy above the majority baseline, is also consistent with partial genuine knowledge. Individual SCORE reports were public before the consolidated paper~\cite{score_g5sny}, so a ``pre-cutoff'' model is not necessarily unexposed.

\emph{Confidence.} Opus 5.5 is the only verifier whose top-decile gain has a CI excluding zero: \llmOpusIdTop{} (CI \llmOpusIdTopCi{}) vs.\ \llmOpusIdAll{}, gain CI \llmOpusIdGainCi{}. This is nominal, since it is unadjusted across at least 11 selective comparisons, and it was obtained on identified input, where memorization would act (de-identified top decile: \llmOpusDeidTop{}). For GPT-6.1 sol the gain CI (\llmSolIdGainCi{}) touches zero, and B1's confident decile is correct only \llmBoneTop{} of the time. Opus 5.5 is thus the strongest candidate for confidence gating, but it is also a judge whose temporal validity cannot be certified.

\subsection{Robustness}
% Criterion-sensitivity table (paper/tables/sens.tex) reported in text for space.
% Full subgroup table: paper/tables/sub.tex (released with the artifact)
Under three alternative success criteria from the SCORE package (analyst judgement, replication within the original CI, Bayesian replication), B1 stays between \sensLo{} and \sensHi{}, and M2 is never consistently better. Within fields, B1 is near chance in psychology and health (AUROC \subPsych{}, $n{=}\subPsychN{}$) and higher in economics (\subEcon{}) and political science (\subPol{}); these small-sample estimates are descriptive. Weighting papers equally, as the SCORE analysis does (one paper contributes \maxClaimsPaper{} claims), gives AUROCs of \pwBone{} (B1), \pwMtwo{} (M2) and \pwDplan{} (D, uninterpreted, no paired CI); graph features remain no better than metadata.

\begin{table}[t]\centering\caption{Agent-in-the-loop triage on identical claims (scenario A: 274 valid queries; B: 259 post-launch queries, where the unprotected tool silently returns a no-cutoff prediction and the contract-aware tool refuses). U: agent sees only \texttt{p\_replicate}; C: full contract output. Wrong auto = accepted non-replicating plus rejected replicating claims, as a share of all claims; $\Delta$ = C$-$U with paper-clustered 95\% CI. $^{\dagger}$Exploratory revised wording (\S\ref{sec:agent}).}\label{tab:agent}
\scriptsize\setlength{\tabcolsep}{2pt}\begin{tabular}{@{}llccccc@{}}\toprule
& Agent & Esc.\ U & Esc.\ C & Wrong U & Wrong C & $\Delta$ wrong (pp) \\\midrule
A & GPT-5.4 mini & 66\% & 93\% & 14.6\% & 3.6\% & -10.9 [-16.4, -6.3] \\
 & GPT-6.1 sol & 81\% & 100\% & 9.1\% & 0.0\% & -9.1 [-13.5, -5.1] \\
 & Claude Haiku 4.5 & 59\% & 1\% & 17.9\% & 54.0\% & +36.1 [+27.3, +44.4] \\
 & Claude Haiku 4.5, C$^{\dagger}$ & & 100\% & & 0.0\% & \\
 & Claude Opus 5.5 & 94\% & 100\% & 2.9\% & 0.0\% & -2.9 [-5.2, -1.1] \\
\midrule
B & GPT-5.4 mini & 54\% & 100\% & 20.1\% & 0.0\% & -20.1 [-29.0, -12.5] \\
 & GPT-6.1 sol & 75\% & 100\% & 12.0\% & 0.0\% & -12.0 [-19.8, -5.6] \\
 & Claude Haiku 4.5 & 49\% & 85\% & 21.6\% & 6.9\% & -14.7 [-22.3, -7.5] \\
 & Claude Haiku 4.5, C$^{\dagger}$ & & 100\% & & 0.0\% & \\
 & Claude Opus 5.5 & 91\% & 100\% & 3.9\% & 0.0\% & -3.9 [-7.5, -0.8] \\
\bottomrule\end{tabular}\end{table}


\section{Agent in the Loop}\label{sec:agent}
Contract tests show that the tool \emph{reports} its limits; they do not show that an agent \emph{uses} them~\cite{hong2026agentactionbench}. We therefore ran a minimal single-turn triage workflow. An LLM agent receives a claim and the output of \texttt{verify\_claim}, and must \emph{accept} the claim (the pipeline builds on it), \emph{reject} it, or \emph{escalate} it to a human reviewer with limited capacity. The system prompt is identical across workflows; only the tool output differs. The \emph{unprotected} agent sees \texttt{p\_replicate} alone; the \emph{contract-aware} agent sees the full response.
In scenario~A, a valid query at the strict cutoff, both agents receive the leave-paper-out prediction for each of the 274 claims. In scenario~B, a query dated 2023-06-01 after the replication launched, the unprotected tool silently returns the no-cutoff prediction, while the contract-aware tool returns its refusal; this covers 259 claims. Scenario~B is deliberately asymmetric: it measures what an agent does when a contract violation is \emph{silent} versus \emph{surfaced}, not harm from leakage itself (the leaky prediction is no more accurate; \S\ref{sec:results}-C). All four LLMs served as agents (\agCalls{} decisions, plus \agCallsVtwo{} in the control, robustness and revised-wording runs below; \$\agCost{} in total).

Table~\ref{tab:agent} shows the results. No verifier probability was extreme (range \agPmin{}--\agPmax{}), so unprotected agents already escalated many claims. Their automatic decisions, however, were near chance and consisted mostly of rejecting claims that did replicate.

For three of the four agents, the contract information turned these errors into escalations: wrong automatic decisions changed by \agMiniADAutoerror{}, \agSolADAutoerror{} and \agOpusADAutoerror{} percentage points in scenario~A (GPT-5.4 mini, GPT-6.1 sol, Opus 5.5), and these agents escalated every post-launch refusal in scenario~B.

\emph{Control arms} (Table~\ref{tab:controls}, scenario~A). Without any tool, agents mostly escalate on their own (GPT-5.4 mini \agMiniANotoolEsc{}, GPT-6.1 sol \agSolANotoolEsc{}, Haiku 4.5 \agHaikuANotoolEsc{}). It is the bare probability that induces automatic decisions: relative to no tool, wrong decisions rose by \agMiniAUnpVsNo{}, \agSolAUnpVsNo{} and \agHaikuAUnpVsNo{} points (CIs \agMiniAUnpVsNoCi{}, \agSolAUnpVsNoCi{}, \agHaikuAUnpVsNoCi{}). A generic ``uncertain'' note undoes part of this harm. The contract does more than that note for three agents, changing wrong decisions relative to the placebo by \agMiniAConVsPla{}, \agSolAConVsPla{} and \agOpusAConVsPla{} points (CIs \agMiniAConVsPlaCi{}, \agSolAConVsPlaCi{}, \agOpusAConVsPlaCi{}; unadjusted). Its effect is therefore not explained by discouraging text alone. But it beats \emph{no tool} only for Opus 5.5 (\agOpusAConVsNo{}, CI \agOpusAConVsNoCi{}), whose no-tool arm accepted some claims; an always-escalate policy matches it. For a near-chance verifier, the contract's value is to remove the false authority of a bare number, not to make automation better. That is a reduction in automatic errors, not an improvement in decision utility: if one escalation costs $c$ relative to one wrong automatic decision, the contract-aware agent beats the unprotected one only when $c$ is below \agMiniABreakeven{}, \agSolABreakeven{} and \agOpusABreakeven{} for these three agents (for Haiku~4.5, never).

\emph{Prompt robustness.} With an independently worded system prompt (three agents, scenario~A), the ordering \emph{bare number} $>$ \emph{placebo} $>$ \emph{contract} in wrong automatic decisions holds for both GPT agents. For GPT-5.4 mini the rates are \agMiniAUPtwoErr{}, \agMiniAPlaPtwoErr{} and \agMiniACPtwoErr{}, and Haiku 4.5's misuse persists: it rejects \agHaikuACPtwoRej{} of claims, with \agHaikuACPtwoErr{} wrong. Haiku's placebo does not help (\agHaikuAPlaPtwoErr{} vs.\ \agHaikuAUPtwoErr{} bare). Magnitudes are prompt-sensitive, however: GPT-6.1 sol's bare-number error rate falls from \agSolAUAutoerror{} to \agSolAUPtwoErr{}.

\begin{table}[t]\centering\caption{Control arms for scenario~A (274 claims): share escalated / share of wrong automatic decisions. \emph{No tool}: the agent sees only the claim. \emph{Placebo}: \textnormal{\texttt{p\_replicate}} plus a generic ``this estimate is uncertain'' note. \emph{Contract}: full \textnormal{\texttt{verify\_claim}} output. \emph{Always escalate} needs no model.}\label{tab:controls}
\scriptsize\setlength{\tabcolsep}{3pt}\begin{tabular}{@{}lcccc@{}}\toprule
Agent & No tool & Placebo & Unprotected & Contract \\\midrule
GPT-5.4 mini & 100\% / 0.0\% & 82\% / 8.8\% & 66\% / 14.6\% & 93\% / 3.6\% \\
GPT-6.1 sol & 100\% / 0.0\% & 95\% / 2.9\% & 81\% / 9.1\% & 100\% / 0.0\% \\
Claude Haiku 4.5 & 98\% / 0.4\% & 63\% / 15.7\% & 59\% / 17.9\% & 1\% / 54.0\% \\
Claude Opus 5.5 & 86\% / 4.0\% & 96\% / 1.8\% & 94\% / 2.9\% & 100\% / 0.0\% \\
\midrule
Always escalate & \multicolumn{4}{c}{100\% / 0.0\%} \\
\bottomrule\end{tabular}\end{table}


The fourth agent, Claude Haiku 4.5, \emph{misused} the contract. Shown the full response in scenario~A, it rejected \agHaikuACRej{} of claims and its wrong automatic decisions rose from \agHaikuAUAutoerror{} to \agHaikuACAutoerror{}. Its stated reasons show two misreadings:
(i)~it treated the verifier's weak skill (``AUROC 0.62'') and the advice ``do not gate actions on this score'' as evidence \emph{against the claim} rather than against the tool;
(ii)~it read the field \texttt{selective\_ok} as a warning about ``selective reporting''.
It also rejected \agHaikuBCRej{} of the refusals in scenario~B, as if a refusal were negative evidence. In an exploratory follow-up designed \emph{after} observing this failure, we renamed the field \texttt{confidence\_gating\_supported} and made the advice action-oriented (``this score cannot justify accepting or rejecting the claim on its own; escalate if a decision is needed''), with a refusal that states it carries no information about the claim. Haiku 4.5 then escalated \agHaikuAVtwoEsc{} of claims in both scenarios, with \agHaikuAVtwoErr{} wrong automatic decisions. (Listing~\ref{lst:tool} shows the original wording, used in all main experiments.) A correct tool is therefore not sufficient: contract fields and advice are themselves an interface that must be tested against the agents that consume them. The experiment is deliberately minimal---two system prompts, single-turn decisions, a single break-even cost model---and the revised wording was tested on one model only.

\section{Discussion: Implications for Agentic Systems}\label{sec:discussion}
\textbf{Machine verifiers are weak signals, and apparently weaker than human crowds.} Our best leakage-safe verifiers reach an AUROC of about 0.63 with a Brier skill below 3\%. On SCORE, human forecasters reached a median AUC of 0.76~\cite{mody2026} (on different evaluation sets, so the comparison is indicative). An orchestrator should not prune hypotheses or literature on such a score; better uses are to prioritise claims for human or experimental checking and to surface the score with its base rate and evidence record. An LLM judge is the tempting replacement but not a drop-in one. The post-cutoff judges ranked best, yet part of their advantage may have depended on recognising the paper, and the cutoff they were told to respect cannot be enforced on their weights; the metamorphic test of \S\ref{sec:tool} has no analogue for knowledge held in parameters. Calibration varied by model, and two models named outcomes for papers that were never replicated, so an agent that asks an LLM ``did this replicate?'' can receive confident, false provenance. At minimum, LLM judges should be wrapped in the same contract: dated prompts, de-identified controls, decoy-calibrated memorization probes, base-rate reporting and audit logs, with recalibration before their probabilities are used.

\textbf{Make temporal validity a type, and test it at the tool boundary.} Most of the engineering effort went into enforcing R1--R3, not into modelling. Availability typing caught the achieved-size leak and the undated TOP Factor fields. Metamorphic tests at the tool boundary then provide regression-tested assurance that the deployed tool honours the cutoff the evaluation assumed. Agent tool descriptions should carry an explicit \texttt{as\_of} parameter and an output schema with the base rate, evidence cutoff, gating verdict and provenance.

\textbf{Add a consumer-side clause to the contract.} The same output shifted three agents to escalation and misled a fourth, which read tool-quality warnings as evidence about the claim. We therefore propose a sixth requirement, \textbf{R6 Consumability}: the output must state the action it licenses (e.g., ``escalate''), not only the actions it forbids, and the contract must be validated against the agents that will consume it. R6 is a \emph{semantic} analogue of consumer-driven contract testing in service engineering~\cite{robinson2006}: there, consumers' expectations of a provider's interface are tested, while here LLM consumers' \emph{interpretation} of the output is tested. Our evidence for R6 comes from a single exploratory revision on one model.

\textbf{Ship negative controls, contract tests and audits with the tool.} The permutation null bounded the graph block's signal; mutation analysis exposed test gaps (and unreliable tool verdicts until re-run); a filesystem audit found two paper identifiers differing only in case (\texttt{RJb}, \texttt{rjb}) that had silently shared one cache entry (fixed throughout). Each check belongs in a verifier's continuous-integration suite, alongside data checksums and a label gate that reconciles to published totals.

\section{Threats to Validity and Limitations}\label{sec:limits}
\textbf{Sample size.} With 274 claims in 164 clusters, CIs are wide, and AUROC differences below roughly 0.04 cannot be resolved. Our nulls rule out large gains, not small ones.
\textbf{Construct.} Replication success depends on the criterion, and a failed replication does not imply a false claim; we report three alternative criteria.
\textbf{Scope.} SCORE samples social and behavioural science papers from 2009--2018 that were selected for replication. Results may not transfer to biomedicine, computational reproducibility or formal verification, which need different labels. Which claims SCORE chose to replicate is itself post-publication information that our cutoff does not control.
\textbf{Evidence.} OpenAlex coverage and dates are imperfect, and author disambiguation errors affect dependency clusters. Our graph lacks semantic support/contradiction and dataset-reuse edges, so H1/H2 are tested only at bibliographic, paper-level granularity.
\textbf{Models.} We evaluated deliberately simple, auditable models. The text model is a basic TF-IDF; richer encoders might extract more signal but bring contamination risk. We did not recalibrate the harness's own verifiers. The LLM judges were run with one prompt and the agents with two, one sample per claim, and each model's lowest reasoning setting, so reasoning settings differ across models. \texttt{gpt-6.1-sol} and \texttt{claude-opus-5-5} are dateless identifiers, and knowledge cutoffs are as stated by the vendors. The revised contract wording was designed after observing one agent's failure and tested on that agent only. LLM results may change with prompts, sampling and model versions, and the probe measures only what a model will \emph{say} it knows. Our intervals come from a paper-clustered bootstrap of seed-averaged out-of-fold predictions. They describe uncertainty over claims and papers, but may not fully reflect model-fitting and hyperparameter-selection variability; selective accuracy uses the same predictions.
\textbf{Tool and agent.} The MCP tool is a research prototype. The agent experiment is a single-turn workflow with two system prompts, evaluated offline. Its controls (no tool, generic warning, always escalate) cover scenario~A only, and the revised wording literally says ``escalate'', so its 100\% escalation partly reflects instruction following. The pre/post-cutoff split is confounded, because individual SCORE reports predate the consolidated paper. The methodology is demonstrated on one tool class; the schema and test relations are tool-agnostic, but applying them to a second tool is future work. H1--H3 were not preregistered, and analyses outside the Holm family are exploratory. Targeted mutants cover the five guarantees we implemented; the standard-tool analysis leaves \mutSurv{} surviving mutants outside the contract's scope (\S\ref{sec:tool}). A mutation score measures how well the tests detect injected faults; it is not proof that the contract holds in every case.

\section{Conclusion}
The central lesson is methodological. A verifier tool for scientific agents should come with an epistemic contract, tests that enforce it at the tool boundary, and evidence that agents read it as intended. Each layer caught something the others missed: availability typing caught a leak in our own pipeline, mutation analysis exposed gaps in our tests (including at the cutoff boundary), and the agent experiment revealed a contract--consumer gap that no tool-level test can see. On the SCORE testbed, a contract-compliant verifier discriminates only modestly, bibliographic evidence graphs yield no statistically established improvement, and confident predictions are not reliable enough to gate agent actions. The best LLM judges rank claims competitively, but their temporal validity cannot be certified, and two of the four judges named outcomes for papers that were never replicated. In the loop, a bare probability induced errors that most agents did not make without the tool, and the contract removed most of them for three of four agents. Replication forecasts should inform scientific agents, not gate them.

\section*{Artifact Availability}
A reproducibility package (code, MCP server, tests, provenance manifest, per-claim predictions, cached LLM and agent responses, and the scripts that generate every table and in-text number) is publicly available at \url{https://github.com/luxinlabs/scireligraph} (MIT licence). SCORE data are available from OSF~\cite{score_g5sny}.

\section*{Acknowledgment}
Portions of the data pipeline, analysis code and manuscript text were drafted with the assistance of an AI system (Claude, Anthropic). The author reviewed and verified all content; all numbers are generated from the released analysis outputs.

\bibliographystyle{IEEEtran}
\bibliography{references}

\end{document}
